\documentclass[aspectratio=169,8pt]{beamer}
\usetheme{Madrid}
\usepackage[utf8]{inputenc}
\usepackage[T1]{fontenc}
\usepackage{amsmath,amssymb}
\usepackage{booktabs}
\usepackage{enumitem}
\setlist[enumerate]{label=\Alph*.,leftmargin=*,itemsep=0pt,topsep=1pt}
\setlist[itemize]{leftmargin=*,itemsep=0pt,topsep=1pt}
\setbeamerfont{block body}{size=\tiny}
\setbeamerfont{block title}{size=\scriptsize}
\setbeamerfont{frametitle}{size=\footnotesize}
\setbeamerfont{normal text}{size=\tiny}
\setbeamertemplate{navigation symbols}{}
\setbeamertemplate{itemize item}{\tiny$\bullet$}
\setbeamertemplate{itemize subitem}{\tiny$\circ$}

\title[GARP RAI]{GARP Risk \& AI --- 300 Practice Questions}
\subtitle{Unofficial Exam Prep}
\author{Unofficial}
\date{\today}

\begin{document}
	\begin{frame}\titlepage\end{frame}
	
	% =================================================================
	\section{Module 1: AI and Risk --- Introduction \& Overview}
	% =================================================================
	
	\begin{frame}{Q1 --- Inscrutability}
		\begin{block}{Question}
			A risk team is concerned that a deep learning fraud model cannot be understood internally, even though it produces accurate predictions. Which term best describes this property?
		\end{block}
		\begin{enumerate}
			\item Explainability
			\item Inscrutability
			\item Interpretability
			\item Generalisability
		\end{enumerate}
		\begin{exampleblock}{Answer: B}\end{exampleblock}
		\begin{block}{Explanation}
			Inscrutability = inability to understand the model's internal mechanisms. Different from explainability (explaining specific decisions) and interpretability (understanding internal logic).
		\end{block}
	\end{frame}
	
	\begin{frame}{Q2 --- Deep Blue}
		\begin{block}{Question}
			In 1997, IBM's Deep Blue defeated Garry Kasparov at chess. Which characteristic of classical AI does this best illustrate?
		\end{block}
		\begin{enumerate}
			\item Deep learning's ability to self-train
			\item Classical AI's reliance on rules and search in well-defined contexts
			\item Gradient descent optimisation
			\item Reinforcement learning from environment
		\end{enumerate}
		\begin{exampleblock}{Answer: B}\end{exampleblock}
		\begin{block}{Explanation}
			Deep Blue was a classical AI system using pre-defined rules and search strategies within the narrow domain of chess.
		\end{block}
	\end{frame}
	
	\begin{frame}{Q3 --- Societal Risk}
		\begin{block}{Question}
			A social media platform's algorithm amplifies misleading political content, affecting public opinion at large. This is best classified as:
		\end{block}
		\begin{enumerate}
			\item Individual risk
			\item Organizational risk
			\item Societal risk
			\item Operational risk
		\end{enumerate}
		\begin{exampleblock}{Answer: C}\end{exampleblock}
		\begin{block}{Explanation}
			Harm to society at large (democratic processes, public discourse) = societal risk. Different from harm to a single user or one company.
		\end{block}
	\end{frame}
	
	\begin{frame}{Q4 --- Deep Learning Breakthrough}
		\begin{block}{Question}
			Which of the following is a distinguishing feature of deep learning compared with classical AI?
		\end{block}
		\begin{enumerate}
			\item Use of A* search algorithm
			\item Use of recursive adversarial search
			\item Self-training via gradient descent on data
			\item Lookahead strategy
		\end{enumerate}
		\begin{exampleblock}{Answer: C}\end{exampleblock}
		\begin{block}{Explanation}
			Deep learning learns from data using backpropagation and gradient descent, unlike classical AI which relies on pre-programmed rules.
		\end{block}
	\end{frame}
	
	\begin{frame}{Q5 --- Data Types}
		\begin{block}{Question}
			A survey asks customers to select the types of investments they hold (stocks, bonds, ETFs). How should these responses be classified?
		\end{block}
		\begin{enumerate}
			\item Ordinal data
			\item Nominal data
			\item Interval data
			\item Ratio data
		\end{enumerate}
		\begin{exampleblock}{Answer: B}\end{exampleblock}
		\begin{block}{Explanation}
			Categories with no inherent ordering = nominal data. Ordinal has ordering; interval and ratio are numeric.
		\end{block}
	\end{frame}
	
	\begin{frame}{Q6 --- Data Scaling}
		\begin{block}{Question}
			Which statement about data scaling is correct?
		\end{block}
		\begin{enumerate}
			\item Standardization creates values bounded by 0 and 1
			\item Normalization transforms data to have mean 0 and SD 1
			\item Standardization gives mean 0 and SD 1; normalization gives [0, 1]
			\item Scaling destroys correlations
		\end{enumerate}
		\begin{exampleblock}{Answer: C}\end{exampleblock}
		\begin{block}{Explanation}
			Standardization: $z=(x-\mu)/\sigma$, mean 0, SD 1. Normalization: min-max to [0, 1]. Both preserve correlations.
		\end{block}
	\end{frame}
	
	\begin{frame}{Q7 --- Stepwise Regression}
		\begin{block}{Question}
			Forward stepwise regression yields AIC values: Step 2: Age+Age$^3$ = 3123.20; Step 3: Age+Age$^3$+Age$^5$ = 3124.50. Which model should be selected?
		\end{block}
		\begin{enumerate}
			\item Age + Age$^3$ + Age$^5$
			\item Age + Age$^3$
			\item Age$^3$ only
			\item Restart with backward elimination
		\end{enumerate}
		\begin{exampleblock}{Answer: B}\end{exampleblock}
		\begin{block}{Explanation}
			Forward stepwise stops when adding a variable no longer reduces AIC. Lower AIC = better model.
		\end{block}
	\end{frame}
	
	\begin{frame}{Q8 --- ML vs Econometrics}
		\begin{block}{Question}
			Which is the primary focus of machine learning compared with classical econometrics?
		\end{block}
		\begin{enumerate}
			\item Inference and parameter significance
			\item Out-of-sample prediction accuracy
			\item Hypothesis testing
			\item Assumption validation
		\end{enumerate}
		\begin{exampleblock}{Answer: B}\end{exampleblock}
		\begin{block}{Explanation}
			ML prioritizes prediction accuracy; econometrics focuses on inference and hypothesis testing.
		\end{block}
	\end{frame}
	
	\begin{frame}{Q9 --- Feature Scaling}
		\begin{block}{Question}
			A KNN model with features age (20–80) and salary (\$20k–\$500k) is dominated by salary. What is the best fix?
		\end{block}
		\begin{enumerate}
			\item Remove the salary feature
			\item Apply standardization or normalization
			\item Use decision trees instead
			\item Reduce the number of observations
		\end{enumerate}
		\begin{exampleblock}{Answer: B}\end{exampleblock}
		\begin{block}{Explanation}
			Distance-based algorithms require scaling. Standardization or normalization prevents any feature from dominating.
		\end{block}
	\end{frame}
	
	\begin{frame}{Q10 --- PCA}
		\begin{block}{Question}
			What is the primary purpose of Principal Component Analysis?
		\end{block}
		\begin{enumerate}
			\item Increase number of features
			\item Reduce dimensionality while preserving variance
			\item Classify observations
			\item Cluster data points
		\end{enumerate}
		\begin{exampleblock}{Answer: B}\end{exampleblock}
		\begin{block}{Explanation}
			PCA creates orthogonal components ordered by variance, allowing dimensionality reduction.
		\end{block}
	\end{frame}
	
	\begin{frame}{Q11 --- Data Cleaning}
		\begin{block}{Question}
			Which is the most common problem encountered in data preparation?
		\end{block}
		\begin{enumerate}
			\item Duplicate observations
			\item Missing data
			\item Outliers
			\item Inconsistent formatting
		\end{enumerate}
		\begin{exampleblock}{Answer: B}\end{exampleblock}
		\begin{block}{Explanation}
			Missing data is the most common problem. Understanding whether missingness is informative is critical.
		\end{block}
	\end{frame}
	
	\begin{frame}{Q12 --- Outliers}
		\begin{block}{Question}
			Which types of models are generally robust to outliers?
		\end{block}
		\begin{enumerate}
			\item Ordinary Least Squares
			\item Tree-based models and SVMs
			\item Polynomial regression
			\item Ridge regression
		\end{enumerate}
		\begin{exampleblock}{Answer: B}\end{exampleblock}
		\begin{block}{Explanation}
			Tree-based models split on thresholds, SVM uses margins. OLS and polynomial regressions are sensitive.
		\end{block}
	\end{frame}
	
	\begin{frame}{Q13 --- Q-Q Plot}
		\begin{block}{Question}
			A Q-Q plot shows tails deviating significantly from the diagonal line. This indicates:
		\end{block}
		\begin{enumerate}
			\item Normality
			\item Fat-tailed distribution
			\item Homoscedasticity
			\item Perfect fit
		\end{enumerate}
		\begin{exampleblock}{Answer: B}\end{exampleblock}
		\begin{block}{Explanation}
			Deviation in tails = fat tails. Financial returns typically show this behaviour.
		\end{block}
	\end{frame}
	
	\begin{frame}{Q14 --- Log Transformation}
		\begin{block}{Question}
			A feature has max/min ratio of 140 and positive skew. What is the most appropriate transformation?
		\end{block}
		\begin{enumerate}
			\item Remove the feature
			\item Take the natural log
			\item Leave as-is
			\item Apply min-max normalization
		\end{enumerate}
		\begin{exampleblock}{Answer: B}\end{exampleblock}
		\begin{block}{Explanation}
			Rule of thumb: max/min > 10 = highly skewed. Log transformation reduces skew.
		\end{block}
	\end{frame}
	
	\begin{frame}{Q15 --- Discretization}
		\begin{block}{Question}
			Converting continuous market capitalization into Small/Medium/Large categories is called:
		\end{block}
		\begin{enumerate}
			\item Standardization
			\item Discretization (binning)
			\item Normalization
			\item One-hot encoding
		\end{enumerate}
		\begin{exampleblock}{Answer: B}\end{exampleblock}
		\begin{block}{Explanation}
			Discretization (binning) converts numeric to categorical variables.
		\end{block}
	\end{frame}
	
	\begin{frame}{Q16 --- One-Hot Encoding}
		\begin{block}{Question}
			Why use one-hot encoding instead of integer codes for nominal categories?
		\end{block}
		\begin{enumerate}
			\item Faster training
			\item Avoids implying false ordering
			\item Fewer features
			\item Handles missing data
		\end{enumerate}
		\begin{exampleblock}{Answer: B}\end{exampleblock}
		\begin{block}{Explanation}
			Integer codes imply ordinality; one-hot avoids this for nominal categories.
		\end{block}
	\end{frame}
	
	\begin{frame}{Q17 --- Training/Validation/Test}
		\begin{block}{Question}
			What is the primary purpose of the validation set?
		\end{block}
		\begin{enumerate}
			\item Estimate model parameters
			\item Select between competing models
			\item Final performance estimate
			\item Data cleaning
		\end{enumerate}
		\begin{exampleblock}{Answer: B}\end{exampleblock}
		\begin{block}{Explanation}
			Training set estimates parameters; validation set selects models; test set gives final unbiased estimate.
		\end{block}
	\end{frame}
	
	\begin{frame}{Q18 --- Multicollinearity}
		\begin{block}{Question}
			Two features have correlation coefficient r = 0.98. What is the main concern?
		\end{block}
		\begin{enumerate}
			\item Independence
			\item Multicollinearity
			\item Non-linearity
			\item Heteroskedasticity
		\end{enumerate}
		\begin{exampleblock}{Answer: B}\end{exampleblock}
		\begin{block}{Explanation}
			High correlation between features inflates standard errors and destabilises coefficient estimates.
		\end{block}
	\end{frame}
	
	\begin{frame}{Q19 --- Longitudinal Data}
		\begin{block}{Question}
			Time-series stock returns are an example of:
		\end{block}
		\begin{enumerate}
			\item Cross-sectional data
			\item Longitudinal data
			\item Nominal data
			\item Categorical data
		\end{enumerate}
		\begin{exampleblock}{Answer: B}\end{exampleblock}
		\begin{block}{Explanation}
			Longitudinal (time-series) data has chronological ordering; cross-sectional is a snapshot.
		\end{block}
	\end{frame}
	
	\begin{frame}{Q20 --- Semi-Structured Data}
		\begin{block}{Question}
			Which is an example of semi-structured data?
		\end{block}
		\begin{enumerate}
			\item Census tables
			\item Emails and HTML files
			\item Free-form text
			\item Sensor data
		\end{enumerate}
		\begin{exampleblock}{Answer: B}\end{exampleblock}
		\begin{block}{Explanation}
			Semi-structured data has some structure (e.g. tags, headers) but not fixed tabular form.
		\end{block}
	\end{frame}
	
	\begin{frame}{Q21 --- PCA Elbow}
		\begin{block}{Question}
			Why retain only components before the elbow in a PCA scree plot?
		\end{block}
		\begin{enumerate}
			\item To maximise the number of components
			\item To capture most variance with fewest components
			\item To remove outliers
			\item To ensure normality
		\end{enumerate}
		\begin{exampleblock}{Answer: B}\end{exampleblock}
		\begin{block}{Explanation}
			Components before the elbow capture most variance; after the elbow, additional components add little information.
		\end{block}
	\end{frame}
	
	\begin{frame}{Q22 --- PCA Loadings}
		\begin{block}{Question}
			PC1 has near-equal positive loadings on all Treasury yields. This component represents:
		\end{block}
		\begin{enumerate}
			\item Slope
			\item Curvature
			\item Level
			\item Noise
		\end{enumerate}
		\begin{exampleblock}{Answer: C}\end{exampleblock}
		\begin{block}{Explanation}
			Equal positive loadings = parallel shift = level. PC2 = slope; PC3 = curvature.
		\end{block}
	\end{frame}
	
	\begin{frame}{Q23 --- Standardization vs PCA}
		\begin{block}{Question}
			Why is it important to standardize data before applying PCA?
		\end{block}
		\begin{enumerate}
			\item To preserve correlations
			\item To prevent features with larger scales from dominating
			\item To avoid overfitting
			\item To improve interpretability
		\end{enumerate}
		\begin{exampleblock}{Answer: B}\end{exampleblock}
		\begin{block}{Explanation}
			PCA summarizes variance. Unscaled features with large magnitude dominate incorrectly.
		\end{block}
	\end{frame}
	
	\begin{frame}{Q24 --- GIGO Principle}
		\begin{block}{Question}
			The acronym GIGO stands for:
		\end{block}
		\begin{enumerate}
			\item Garbage In, Garbage Out
			\item Great Input, Great Output
			\item Generic Input, Generic Output
			\item Gradient In, Gradient Out
		\end{enumerate}
		\begin{exampleblock}{Answer: A}\end{exampleblock}
		\begin{block}{Explanation}
			GIGO emphasizes that model quality is only as good as the input data.
		\end{block}
	\end{frame}
	
	\begin{frame}{Q25 --- Data Preparation}
		\begin{block}{Question}
			Which step is NOT part of data cleaning?
		\end{block}
		\begin{enumerate}
			\item Removing duplicates
			\item Handling missing values
			\item Checking outliers
			\item Applying PCA
		\end{enumerate}
		\begin{exampleblock}{Answer: D}\end{exampleblock}
		\begin{block}{Explanation}
			PCA is a dimensionality reduction technique, not a data cleaning step.
		\end{block}
	\end{frame}
	
	% =================================================================
	\section{Module 2: Tools \& Techniques}
	% =================================================================
	
	\begin{frame}{Q26 --- Optimal K}
		\begin{block}{Question}
			WCSS values: K=1: 1000; K=2: 400; K=3: 380; K=4: 375. What is the optimal K?
		\end{block}
		\begin{enumerate}
			\item 1
			\item 2
			\item 3
			\item 4
		\end{enumerate}
		\begin{exampleblock}{Answer: B}\end{exampleblock}
		\begin{block}{Explanation}
			The elbow is at K=2. Further reductions in WCSS are marginal relative to added complexity.
		\end{block}
	\end{frame}
	
	\begin{frame}{Q27 --- Local Optima}
		\begin{block}{Question}
			Running K-means 20 times yields WCSS: 98, 100, 101, 102, 250. Which statement is correct?
		\end{block}
		\begin{enumerate}
			\item 250 is best because it explored a different region
			\item 250 likely converged to a poor local optimum
			\item All are equivalent
			\item 250 is impossible
		\end{enumerate}
		\begin{exampleblock}{Answer: B}\end{exampleblock}
		\begin{block}{Explanation}
			K-means can converge to local optima. Choose lowest WCSS. K-means++ mitigates this.
		\end{block}
	\end{frame}
	
	\begin{frame}{Q28 --- WCSS and BCSS}
		\begin{block}{Question}
			Which statement about WCSS and BCSS is correct?
		\end{block}
		\begin{enumerate}
			\item Both should be minimised
			\item WCSS = separation; BCSS = compactness
			\item Lower WCSS and higher BCSS indicate better clustering
			\item They are always equal
		\end{enumerate}
		\begin{exampleblock}{Answer: C}\end{exampleblock}
		\begin{block}{Explanation}
			WCSS = compactness (lower better). BCSS = separation (higher better).
		\end{block}
	\end{frame}
	
	\begin{frame}{Q29 --- Dendrogram}
		\begin{block}{Question}
			A long vertical line near the top of a dendrogram indicates:
		\end{block}
		\begin{enumerate}
			\item No cluster structure
			\item Merging would significantly worsen fit
			\item Clusters are spherical
			\item K-means gives identical results
		\end{enumerate}
		\begin{exampleblock}{Answer: B}\end{exampleblock}
		\begin{block}{Explanation}
			Long vertical lines show significant cluster boundaries; cutting above yields the clearest clusters.
		\end{block}
	\end{frame}
	
	\begin{frame}{Q30 --- DBSCAN}
		\begin{block}{Question}
			Very small $\varepsilon$ and low minPts in DBSCAN most likely causes:
		\end{block}
		\begin{enumerate}
			\item Too many core points
			\item Fragmented clusters and many noise points
			\item Spherical clusters
			\item Faster convergence
		\end{enumerate}
		\begin{exampleblock}{Answer: B}\end{exampleblock}
		\begin{block}{Explanation}
			Small radius fragments clusters; many points become noise.
		\end{block}
	\end{frame}
	
	\begin{frame}{Q31 --- Concentric Circles}
		\begin{block}{Question}
			Two concentric circles of data. Which algorithm is best suited to find the two clusters?
		\end{block}
		\begin{enumerate}
			\item K-means (Euclidean)
			\item K-means (Manhattan)
			\item DBSCAN
			\item Linear discriminant analysis
		\end{enumerate}
		\begin{exampleblock}{Answer: C}\end{exampleblock}
		\begin{block}{Explanation}
			K-means assumes spherical clusters. DBSCAN uses density and handles arbitrary shapes.
		\end{block}
	\end{frame}
	
	\begin{frame}{Q32 --- Centroid}
		\begin{block}{Question}
			What is the centroid of the points (1, 3), (2, 4), (3, 5)?
		\end{block}
		\begin{enumerate}
			\item (2, 4)
			\item (1.5, 3.5)
			\item (2, 5)
			\item (3, 4)
		\end{enumerate}
		\begin{exampleblock}{Answer: A}\end{exampleblock}
		\begin{block}{Explanation}
			Mean of coordinates: (1+2+3)/3 = 2; (3+4+5)/3 = 4.
		\end{block}
	\end{frame}
	
	\begin{frame}{Q33 --- Silhouette Score}
		\begin{block}{Question}
			A point has $a_i = 2$ (distance to own cluster) and $b_i = 8$ (distance to nearest other cluster). Its silhouette score is:
		\end{block}
		\begin{enumerate}
			\item $-0.75$
			\item $0.25$
			\item $0.75$
			\item $1.00$
		\end{enumerate}
		\begin{exampleblock}{Answer: C}\end{exampleblock}
		\begin{block}{Explanation}
			$s=(b-a)/\max(a,b)=(8-2)/8=0.75$. Values near 1 indicate well-clustered points.
		\end{block}
	\end{frame}
	
	\begin{frame}{Q34 --- Convolution Output}
		\begin{block}{Question}
			A 5×5 input is convolved with a 3×3 kernel (stride 1, no padding). What is the output size?
		\end{block}
		\begin{enumerate}
			\item 5×5
			\item 4×4
			\item 3×3
			\item 2×2
		\end{enumerate}
		\begin{exampleblock}{Answer: C}\end{exampleblock}
		\begin{block}{Explanation}
			$(n-k+1)\times(n-k+1)=(5-3+1)\times(5-3+1)=3\times 3$.
		\end{block}
	\end{frame}
	
	\begin{frame}{Q35 --- Bagging vs Random Forest}
		\begin{block}{Question}
			100 trees are trained on bootstrap samples using all features at every split. This is:
		\end{block}
		\begin{enumerate}
			\item Random forest
			\item Bagging
			\item Gradient boosting
			\item AdaBoost
		\end{enumerate}
		\begin{exampleblock}{Answer: B}\end{exampleblock}
		\begin{block}{Explanation}
			Random forest adds random feature subsetting. Using all features = bagging only.
		\end{block}
	\end{frame}
	
	\begin{frame}{Q36 --- Random Forest Advantage}
		\begin{block}{Question}
			Why does random forest typically outperform a single decision tree?
		\end{block}
		\begin{enumerate}
			\item It has fewer hyperparameters
			\item It reduces correlation across trees
			\item It guarantees higher training accuracy
			\item It eliminates overfitting entirely
		\end{enumerate}
		\begin{exampleblock}{Answer: B}\end{exampleblock}
		\begin{block}{Explanation}
			Random feature subsetting decorrelates trees, reducing variance and improving generalisation.
		\end{block}
	\end{frame}
	
	\begin{frame}{Q37 --- SVM Classification}
		\begin{block}{Question}
			Boundary: $0.5x_1 - 1.2x_2 + 0.3 = 0$. For point $(x_1=2, x_2=1)$, the correct classification is:
		\end{block}
		\begin{enumerate}
			\item Class 1
			\item Class 0
			\item Cannot determine
			\item Boundary case
		\end{enumerate}
		\begin{exampleblock}{Answer: A}\end{exampleblock}
		\begin{block}{Explanation}
			$0.5(2) - 1.2(1) + 0.3 = 1 - 1.2 + 0.3 = 0.1 > 0$ → Class 1.
		\end{block}
	\end{frame}
	
	\begin{frame}{Q38 --- Precision and Recall}
		\begin{block}{Question}
			Precision = 0.9, Recall = 0.4. This means:
		\end{block}
		\begin{enumerate}
			\item 90\% of actual frauds caught; 40\% alerts real
			\item 90\% of alerts are real frauds; 40\% of actual frauds caught
			\item 90\% accuracy
			\item 40\% alerts real; 90\% frauds caught
		\end{enumerate}
		\begin{exampleblock}{Answer: B}\end{exampleblock}
		\begin{block}{Explanation}
			Precision = TP/(TP+FP); Recall = TP/(TP+FN). High precision = few false alerts; low recall = many missed frauds.
		\end{block}
	\end{frame}
	
	\begin{frame}{Q39 --- F1 Score}
		\begin{block}{Question}
			Model A: precision = 0.5, recall = 0.5. Model B: precision = 0.9, recall = 0.1. Which has higher F1?
		\end{block}
		\begin{enumerate}
			\item Model A
			\item Model B
			\item Equal
			\item Cannot compute
		\end{enumerate}
		\begin{exampleblock}{Answer: A}\end{exampleblock}
		\begin{block}{Explanation}
			F1 = 2PR/(P+R). A: 0.5. B: 0.18. A is higher.
		\end{block}
	\end{frame}
	
	\begin{frame}{Q40 --- AUC}
		\begin{block}{Question}
			Model A has AUC = 0.85; Model B has AUC = 0.55. Which statement is correct?
		\end{block}
		\begin{enumerate}
			\item Model B is only slightly better than random
			\item Model B has better discrimination
			\item AUC = 0.55 indicates perfect classification
			\item AUC measures regression accuracy
		\end{enumerate}
		\begin{exampleblock}{Answer: A}\end{exampleblock}
		\begin{block}{Explanation}
			AUC = 0.5 = random. Model A clearly outperforms Model B.
		\end{block}
	\end{frame}
	
	\begin{frame}{Q41 --- Confusion Matrix Accuracy}
		\begin{block}{Question}
			Confusion matrix: TP=40, FP=10, FN=20, TN=130. What is the accuracy?
		\end{block}
		\begin{enumerate}
			\item 0.80
			\item 0.85
			\item 0.90
			\item 0.95
		\end{enumerate}
		\begin{exampleblock}{Answer: B}\end{exampleblock}
		\begin{block}{Explanation}
			Accuracy = $(40+130)/(40+10+20+130) = 170/200 = 0.85$.
		\end{block}
	\end{frame}
	
	\begin{frame}{Q42 --- Tokenization}
		\begin{block}{Question}
			Which NLP pre-processing step splits text into individual words?
		\end{block}
		\begin{enumerate}
			\item Lemmatization
			\item Stop word removal
			\item Tokenization
			\item Feature extraction
		\end{enumerate}
		\begin{exampleblock}{Answer: C}\end{exampleblock}
		\begin{block}{Explanation}
			Tokenization is the first step: splits text into words/tokens.
		\end{block}
	\end{frame}
	
	\begin{frame}{Q43 --- Skip-gram}
		\begin{block}{Question}
			Which Word2Vec architecture has one input neuron and multiple output neurons?
		\end{block}
		\begin{enumerate}
			\item CBoW
			\item Skip-gram
			\item Autoencoder
			\item RNN
		\end{enumerate}
		\begin{exampleblock}{Answer: B}\end{exampleblock}
		\begin{block}{Explanation}
			CBoW: many inputs, one output. Skip-gram: one input, many outputs.
		\end{block}
	\end{frame}
	
	\begin{frame}{Q44 --- Co-training}
		\begin{block}{Question}
			What is the main benefit of co-training?
		\end{block}
		\begin{enumerate}
			\item Faster training via parallel computing
			\item Exploits unlabelled data using two feature views
			\item Reduces feature count
			\item Eliminates the need for validation
		\end{enumerate}
		\begin{exampleblock}{Answer: B}\end{exampleblock}
		\begin{block}{Explanation}
			Two models on different feature views label data for each other, exploiting unlabelled data.
		\end{block}
	\end{frame}
	
	\begin{frame}{Q45 --- Ridge vs LASSO}
		\begin{block}{Question}
			A model has 100 features, only 10 of which matter. Which regularisation is best?
		\end{block}
		\begin{enumerate}
			\item Ridge (L2)
			\item LASSO (L1)
			\item Neither
			\item Both are identical
		\end{enumerate}
		\begin{exampleblock}{Answer: B}\end{exampleblock}
		\begin{block}{Explanation}
			LASSO shrinks irrelevant coefficients to exactly zero (feature selection). Ridge only shrinks toward zero.
		\end{block}
	\end{frame}
	
	\begin{frame}{Q46 --- Elastic Net}
		\begin{block}{Question}
			A dataset has 50 correlated features plus 10 truly important features. Which regularisation is best?
		\end{block}
		\begin{enumerate}
			\item Ridge only
			\item LASSO only
			\item Elastic Net
			\item No regularisation
		\end{enumerate}
		\begin{exampleblock}{Answer: C}\end{exampleblock}
		\begin{block}{Explanation}
			Elastic Net combines L1 (feature selection) and L2 (handles correlated features).
		\end{block}
	\end{frame}
	
	\begin{frame}{Q47 --- Cross-Validation}
		\begin{block}{Question}
			With 100 observations and k = 10 folds, how many observations are in each validation fold?
		\end{block}
		\begin{enumerate}
			\item 1
			\item 5
			\item 10
			\item 20
		\end{enumerate}
		\begin{exampleblock}{Answer: C}\end{exampleblock}
		\begin{block}{Explanation}
			100/10 = 10 observations per fold.
		\end{block}
	\end{frame}
	
	\begin{frame}{Q48 --- Stratified CV}
		\begin{block}{Question}
			A dataset is 95\% non-default and 5\% default. Why is stratified cross-validation important?
		\end{block}
		\begin{enumerate}
			\item It speeds up training
			\item It ensures each fold contains proportional class representation
			\item It removes outliers
			\item It eliminates overfitting
		\end{enumerate}
		\begin{exampleblock}{Answer: B}\end{exampleblock}
		\begin{block}{Explanation}
			Standard k-fold may leave some folds with zero minority-class examples. Stratification preserves the class ratio.
		\end{block}
	\end{frame}
	
	\begin{frame}{Q49 --- Bias-Variance}
		\begin{block}{Question}
			A model has low training error but high test error. This indicates:
		\end{block}
		\begin{enumerate}
			\item High bias, low variance
			\item High variance, low bias
			\item Optimal fit
			\item Underfitting
		\end{enumerate}
		\begin{exampleblock}{Answer: B}\end{exampleblock}
		\begin{block}{Explanation}
			Low training error + high test error = overfitting = high variance, low bias.
		\end{block}
	\end{frame}
	
	\begin{frame}{Q50 --- Underfitting}
		\begin{block}{Question}
			A linear model is fit to data with a strong U-shaped relationship. This is an example of:
		\end{block}
		\begin{enumerate}
			\item Overfitting
			\item Underfitting
			\item Regularisation
			\item Stratification
		\end{enumerate}
		\begin{exampleblock}{Answer: B}\end{exampleblock}
		\begin{block}{Explanation}
			Linear model cannot capture nonlinearity = underfitting (high bias).
		\end{block}
	\end{frame}
	
	\begin{frame}{Q51 --- Cancer Screening}
		\begin{block}{Question}
			In cancer screening, which metric should be prioritised?
		\end{block}
		\begin{enumerate}
			\item Precision (avoid false alarms)
			\item Recall/Sensitivity (avoid missing cases)
			\item Accuracy (overall)
			\item Specificity
		\end{enumerate}
		\begin{exampleblock}{Answer: B}\end{exampleblock}
		\begin{block}{Explanation}
			Missing a cancer case (false negative) is far more costly than a false alarm. Maximise recall.
		\end{block}
	\end{frame}
	
	\begin{frame}{Q52 --- ROC Curve}
		\begin{block}{Question}
			Which statement about the ROC curve is correct?
		\end{block}
		\begin{enumerate}
			\item Each point corresponds to a different classification threshold
			\item ROC plots training vs test error
			\item Lower AUC is better
			\item ROC is only for regression
		\end{enumerate}
		\begin{exampleblock}{Answer: A}\end{exampleblock}
		\begin{block}{Explanation}
			ROC plots TPR vs FPR across all thresholds. Higher AUC = better discrimination.
		\end{block}
	\end{frame}
	
	\begin{frame}{Q53 --- Specificity}
		\begin{block}{Question}
			Sensitivity = 0.8, Specificity = 0.9. What is the false positive rate?
		\end{block}
		\begin{enumerate}
			\item 0.10
			\item 0.20
			\item 0.80
			\item 0.90
		\end{enumerate}
		\begin{exampleblock}{Answer: A}\end{exampleblock}
		\begin{block}{Explanation}
			FPR = 1 $-$ Specificity = 1 $-$ 0.9 = 0.1.
		\end{block}
	\end{frame}
	
	\begin{frame}{Q54 --- AUC Edge Cases}
		\begin{block}{Question}
			A model has AUC = 0.45. This indicates:
		\end{block}
		\begin{enumerate}
			\item Better than random
			\item Worse than random
			\item Perfect classifier
			\item Random classifier
		\end{enumerate}
		\begin{exampleblock}{Answer: B}\end{exampleblock}
		\begin{block}{Explanation}
			AUC = 0.5 = random. AUC < 0.5 means the model performs worse than random.
		\end{block}
	\end{frame}
	
	\begin{frame}{Q55 --- Threshold Choice}
		\begin{block}{Question}
			A bank lowers the fraud-score threshold. What happens?
		\end{block}
		\begin{enumerate}
			\item Recall increases, precision typically decreases
			\item Precision increases, recall decreases
			\item Both increase
			\item Both decrease
		\end{enumerate}
		\begin{exampleblock}{Answer: A}\end{exampleblock}
		\begin{block}{Explanation}
			Lower threshold = more positives predicted = catches more frauds (higher recall) but more false alarms (lower precision).
		\end{block}
	\end{frame}
	
	\begin{frame}{Q56 --- Regularization Purpose}
		\begin{block}{Question}
			What is the primary purpose of regularisation in machine learning?
		\end{block}
		\begin{enumerate}
			\item To increase model complexity
			\item To prevent overfitting by penalizing large coefficients
			\item To speed up training
			\item To increase the number of features
		\end{enumerate}
		\begin{exampleblock}{Answer: B}\end{exampleblock}
		\begin{block}{Explanation}
			Regularization (Ridge, LASSO, Elastic Net) prevents overfitting by penalizing large coefficients.
		\end{block}
	\end{frame}
	
	\begin{frame}{Q57 --- L1 Penalty}
		\begin{block}{Question}
			Which regularisation technique uses the L1 penalty?
		\end{block}
		\begin{enumerate}
			\item Ridge
			\item LASSO
			\item Both
			\item Neither
		\end{enumerate}
		\begin{exampleblock}{Answer: B}\end{exampleblock}
		\begin{block}{Explanation}
			Ridge = L2 (squared). LASSO = L1 (absolute).
		\end{block}
	\end{frame}
	
	\begin{frame}{Q58 --- Grid Search}
		\begin{block}{Question}
			Grid search is used for:
		\end{block}
		\begin{enumerate}
			\item Feature selection
			\item Hyperparameter tuning
			\item Cross-validation
			\item Data cleaning
		\end{enumerate}
		\begin{exampleblock}{Answer: B}\end{exampleblock}
		\begin{block}{Explanation}
			Grid search exhaustively searches hyperparameter combinations, often with CV.
		\end{block}
	\end{frame}
	
	\begin{frame}{Q59 --- Bootstrapping}
		\begin{block}{Question}
			Sampling with replacement to estimate model performance is called:
		\end{block}
		\begin{enumerate}
			\item Cross-validation
			\item Bootstrapping
			\item Grid search
			\item Stratification
		\end{enumerate}
		\begin{exampleblock}{Answer: B}\end{exampleblock}
		\begin{block}{Explanation}
			Bootstrapping resamples with replacement. About 37\% of data is out-of-bootstrap each iteration.
		\end{block}
	\end{frame}
	
	\begin{frame}{Q60 --- Learning Rate}
		\begin{block}{Question}
			Too-large learning rate in gradient descent most likely causes:
		\end{block}
		\begin{enumerate}
			\item Slow convergence
			\item Overshooting and oscillation
			\item Vanishing gradient
			\item Perfect convergence
		\end{enumerate}
		\begin{exampleblock}{Answer: B}\end{exampleblock}
		\begin{block}{Explanation}
			Too-large n overshoots the minimum, causing oscillation or divergence.
		\end{block}
	\end{frame}
	
	\begin{frame}{Q61 --- Vanishing Gradient}
		\begin{block}{Question}
			Vanishing gradient is most likely in:
		\end{block}
		\begin{enumerate}
			\item Shallow networks
			\item Deep networks with sigmoid activations
			\item Linear regression
			\item Decision trees
		\end{enumerate}
		\begin{exampleblock}{Answer: B}\end{exampleblock}
		\begin{block}{Explanation}
			Chain rule of small derivatives → vanishing gradient. Use ReLU, batch norm, or LSTM to fix.
		\end{block}
	\end{frame}
	
	\begin{frame}{Q62 --- Backpropagation}
		\begin{block}{Question}
			Backpropagation uses which calculus rule?
		\end{block}
		\begin{enumerate}
			\item Bayes rule
			\item Chain rule
			\item Product rule
			\item Power rule
		\end{enumerate}
		\begin{exampleblock}{Answer: B}\end{exampleblock}
		\begin{block}{Explanation}
			Backpropagation applies the chain rule to compute gradients layer by layer.
		\end{block}
	\end{frame}
	
	\begin{frame}{Q63 --- Momentum}
		\begin{block}{Question}
			Momentum helps gradient descent by:
		\end{block}
		\begin{enumerate}
			\item Reducing learning rate
			\item Accelerating in consistent directions and reducing overshoot
			\item Eliminating local minima
			\item Reducing feature count
		\end{enumerate}
		\begin{exampleblock}{Answer: B}\end{exampleblock}
		\begin{block}{Explanation}
			Momentum smooths the descent, accelerating consistent directions and reducing oscillation.
		\end{block}
	\end{frame}
	
	\begin{frame}{Q64 --- ReLU}
		\begin{block}{Question}
			ReLU is preferred over sigmoid in hidden layers because:
		\end{block}
		\begin{enumerate}
			\item It is nonlinear
			\item It reduces vanishing gradient
			\item It is bounded
			\item It is linear
		\end{enumerate}
		\begin{exampleblock}{Answer: B}\end{exampleblock}
		\begin{block}{Explanation}
			ReLU does not saturate for positive values, reducing vanishing gradient.
		\end{block}
	\end{frame}
	
	\begin{frame}{Q65 --- MSE vs MAE}
		\begin{block}{Question}
			Which metric is more sensitive to outliers?
		\end{block}
		\begin{enumerate}
			\item MSE
			\item MAE
			\item MAPE
			\item Both equal
		\end{enumerate}
		\begin{exampleblock}{Answer: A}\end{exampleblock}
		\begin{block}{Explanation}
			MSE squares errors, heavily penalising large deviations.
		\end{block}
	\end{frame}
	
	\begin{frame}{Q66 --- MAPE}
		\begin{block}{Question}
			MAPE fails when:
		\end{block}
		\begin{enumerate}
			\item Data is large
			\item Actual values near zero
			\item Data is skewed
			\item Outliers present
		\end{enumerate}
		\begin{exampleblock}{Answer: B}\end{exampleblock}
		\begin{block}{Explanation}
			Division by near-zero actuals makes MAPE explode or undefined.
		\end{block}
	\end{frame}
	
	\begin{frame}{Q67 --- RMSE Units}
		\begin{block}{Question}
			RMSE is preferred over MSE because:
		\end{block}
		\begin{enumerate}
			\item Lower value
			\item Same units as target
			\item Less sensitive
			\item Easy to compute
		\end{enumerate}
		\begin{exampleblock}{Answer: B}\end{exampleblock}
		\begin{block}{Explanation}
			RMSE is in the same units as the target, making it more interpretable.
		\end{block}
	\end{frame}
	
	\begin{frame}{Q68 --- False Positive}
		\begin{block}{Question}
			A false positive is:
		\end{block}
		\begin{enumerate}
			\item Predicted positive, actually positive
			\item Predicted positive, actually negative
			\item Predicted negative, actually positive
			\item Predicted negative, actually negative
		\end{enumerate}
		\begin{exampleblock}{Answer: B}\end{exampleblock}
		\begin{block}{Explanation}
			FP = Type I error. Predicted positive but actually negative.
		\end{block}
	\end{frame}
	
	\begin{frame}{Q69 --- Specificity}
		\begin{block}{Question}
			Specificity = TN/(TN+FP) measures:
		\end{block}
		\begin{enumerate}
			\item True positive rate
			\item True negative rate
			\item Precision
			\item Accuracy
		\end{enumerate}
		\begin{exampleblock}{Answer: B}\end{exampleblock}
		\begin{block}{Explanation}
			Specificity is the true negative rate.
		\end{block}
	\end{frame}
	
	\begin{frame}{Q70 --- Class Imbalance}
		\begin{block}{Question}
			98\% majority class. A model predicting all majority gets:
		\end{block}
		\begin{enumerate}
			\item Low accuracy
			\item 98\% accuracy but useless
			\item Perfect recall
			\item Poor precision
		\end{enumerate}
		\begin{exampleblock}{Answer: B}\end{exampleblock}
		\begin{block}{Explanation}
			Accuracy paradox: high accuracy but no predictive value.
		\end{block}
	\end{frame}
	
	\begin{frame}{Q71 --- AUC and Imbalance}
		\begin{block}{Question}
			AUC's main advantage over accuracy in imbalanced data:
		\end{block}
		\begin{enumerate}
			\item Faster
			\item Threshold-independent and imbalance-robust
			\item Simpler
			\item Requires fewer features
		\end{enumerate}
		\begin{exampleblock}{Answer: B}\end{exampleblock}
		\begin{block}{Explanation}
			AUC evaluates discrimination across thresholds, unaffected by class ratio.
		\end{block}
	\end{frame}
	
	\begin{frame}{Q72 --- ROC Axes}
		\begin{block}{Question}
			The ROC curve plots:
		\end{block}
		\begin{enumerate}
			\item TPR vs FPR
			\item Precision vs Recall
			\item Error vs Complexity
			\item Train vs Test error
		\end{enumerate}
		\begin{exampleblock}{Answer: A}\end{exampleblock}
		\begin{block}{Explanation}
			True positive rate vs false positive rate across thresholds.
		\end{block}
	\end{frame}
	
	\begin{frame}{Q73 --- Cost Matrix}
		\begin{block}{Question}
			The classification threshold should be set to:
		\end{block}
		\begin{enumerate}
			\item Always 0.5
			\item Minimise expected cost using a cost matrix
			\item Maximise accuracy
			\item Minimise training error
		\end{enumerate}
		\begin{exampleblock}{Answer: B}\end{exampleblock}
		\begin{block}{Explanation}
			Threshold should reflect business costs of FP vs FN.
		\end{block}
	\end{frame}
	
	\begin{frame}{Q74 --- Multi-Class Metrics}
		\begin{block}{Question}
			For 3-class classification, precision is computed by:
		\end{block}
		\begin{enumerate}
			\item Averaging binary precisions
			\item Summing all TP / summing all (TP+FP)
			\item Using majority class only
			\item Ignoring minority classes
		\end{enumerate}
		\begin{exampleblock}{Answer: B}\end{exampleblock}
		\begin{block}{Explanation}
			Micro-average precision sums TP over sums of TP+FP across classes.
		\end{block}
	\end{frame}
	
	\begin{frame}{Q75 --- Polynomial Overfitting}
		\begin{block}{Question}
			Adding too many polynomial terms likely causes:
		\end{block}
		\begin{enumerate}
			\item Underfitting
			\item Overfitting
			\item Regularisation
			\item Stratification
		\end{enumerate}
		\begin{exampleblock}{Answer: B}\end{exampleblock}
		\begin{block}{Explanation}
			High-degree polynomial fits noise = overfitting = high variance.
		\end{block}
	\end{frame}
	
	% =================================================================
	\section{Module 3: Risk and Risk Factors}
	% =================================================================
	
	\begin{frame}{Q76 --- Sensitivity Priority}
		\begin{block}{Question}
			A rare fatal disease affects 0.1\% of the population. Which metric should be prioritised?
		\end{block}
		\begin{enumerate}
			\item Specificity
			\item Precision
			\item Sensitivity
			\item Accuracy
		\end{enumerate}
		\begin{exampleblock}{Answer: C}\end{exampleblock}
		\begin{block}{Explanation}
			False negatives are fatal → maximise true positive rate (sensitivity).
		\end{block}
	\end{frame}
	
	\begin{frame}{Q77 --- Data Composition Bias}
		\begin{block}{Question}
			A fraud model uses 10,000 transactions with only 20 fraudulent. Which bias is most likely?
		\end{block}
		\begin{enumerate}
			\item Measurement bias
			\item Representation bias
			\item Data composition bias
			\item Historical bias
		\end{enumerate}
		\begin{exampleblock}{Answer: C}\end{exampleblock}
		\begin{block}{Explanation}
			Severe class imbalance = data composition bias. Model underperforms on minority class.
		\end{block}
	\end{frame}
	
	\begin{frame}{Q78 --- Fairness Fix}
		\begin{block}{Question}
			A logistic hiring model uses education unfairly. Which remedy ensures equal treatment without removing the variable?
		\end{block}
		\begin{enumerate}
			\item Remove the education level variable
			\item Replace education with its average at prediction time
			\item Add interaction terms
			\item Create dummy variables
		\end{enumerate}
		\begin{exampleblock}{Answer: B}\end{exampleblock}
		\begin{block}{Explanation}
			Replacing with the average at prediction time treats identical candidates equally.
		\end{block}
	\end{frame}
	
	\begin{frame}{Q79 --- Demographic Parity}
		\begin{block}{Question}
			Group A has 30\% loan approvals; Group B has 30\%. Which fairness metric is satisfied?
		\end{block}
		\begin{enumerate}
			\item Equal opportunity
			\item Predictive rate parity
			\item Demographic parity
			\item Individual fairness
		\end{enumerate}
		\begin{exampleblock}{Answer: C}\end{exampleblock}
		\begin{block}{Explanation}
			Demographic parity = equal positive prediction rates across groups.
		\end{block}
	\end{frame}
	
	\begin{frame}{Q80 --- Human-in-the-Loop}
		\begin{block}{Question}
			A chatbot gives incorrect financial advice. Which action best mitigates the risk?
		\end{block}
		\begin{enumerate}
			\item Increase training data
			\item Add human verification before responses are sent
			\item Deploy more frequently
			\item Disable the chatbot
		\end{enumerate}
		\begin{exampleblock}{Answer: B}\end{exampleblock}
		\begin{block}{Explanation}
			Human-in-the-loop catches errors before they harm customers.
		\end{block}
	\end{frame}
	
	\begin{frame}{Q81 --- Model Risk vs Bias}
		\begin{block}{Question}
			Which statement correctly distinguishes model risk from algorithmic bias?
		\end{block}
		\begin{enumerate}
			\item Model risk = errors/misuse; bias = unfair discrimination
			\item Model risk = unfair discrimination; bias = errors
			\item They are identical
			\item Model risk applies only to AI
		\end{enumerate}
		\begin{exampleblock}{Answer: A}\end{exampleblock}
		\begin{block}{Explanation}
			Model risk covers all adverse outcomes from model errors/misuse. Algorithmic bias is unfair discrimination.
		\end{block}
	\end{frame}
	
	\begin{frame}{Q82 --- Explainability}
		\begin{block}{Question}
			Which is the primary reason regulators demand model explainability?
		\end{block}
		\begin{enumerate}
			\item To improve accuracy
			\item To increase speed
			\item For trust, accountability, and compliance
			\item To reduce features
		\end{enumerate}
		\begin{exampleblock}{Answer: C}\end{exampleblock}
		\begin{block}{Explanation}
			Explainability supports trust, accountability, and regulatory compliance (GDPR).
		\end{block}
	\end{frame}
	
	\begin{frame}{Q83 --- Interpretability vs Explainability}
		\begin{block}{Question}
			Interpretability is best defined as:
		\end{block}
		\begin{enumerate}
			\item Explaining specific decisions
			\item Understanding internal model mechanisms
			\item Reducing features
			\item Speeding up training
		\end{enumerate}
		\begin{exampleblock}{Answer: B}\end{exampleblock}
		\begin{block}{Explanation}
			Interpretability = internal mechanisms. Explainability = decisions in human terms.
		\end{block}
	\end{frame}
	
	\begin{frame}{Q84 --- Black Box}
		\begin{block}{Question}
			The black-box problem refers to:
		\end{block}
		\begin{enumerate}
			\item Encrypted model weights
			\item Decisions that cannot be understood by humans
			\item Too-large models for memory
			\item Illegal training data
		\end{enumerate}
		\begin{exampleblock}{Answer: B}\end{exampleblock}
		\begin{block}{Explanation}
			Black-box = opaque internal workings, difficult to interpret.
		\end{block}
	\end{frame}
	
	\begin{frame}{Q85 --- AI in Hiring}
		\begin{block}{Question}
			What is the primary concern with AI in hiring decisions?
		\end{block}
		\begin{enumerate}
			\item AI is too slow
			\item AI may perpetuate historical bias
			\item AI is too expensive
			\item AI cannot process resumes
		\end{enumerate}
		\begin{exampleblock}{Answer: B}\end{exampleblock}
		\begin{block}{Explanation}
			AI trained on historical data may perpetuate past discrimination.
		\end{block}
	\end{frame}
	
	\begin{frame}{Q86 --- Historical Bias}
		\begin{block}{Question}
			AI trained on past hiring data discriminating against women exemplifies:
		\end{block}
		\begin{enumerate}
			\item Measurement bias
			\item Historical bias
			\item Sampling bias
			\item Deployment bias
		\end{enumerate}
		\begin{exampleblock}{Answer: B}\end{exampleblock}
		\begin{block}{Explanation}
			Historical bias = past discrimination encoded in data.
		\end{block}
	\end{frame}
	
	\begin{frame}{Q87 --- Sampling Bias}
		\begin{block}{Question}
			A model trained on US travelers performs poorly in Brazil, Nigeria, and Thailand. Most likely cause:
		\end{block}
		\begin{enumerate}
			\item Problem specification bias
			\item Deployment bias
			\item Historical bias
			\item Sampling bias
		\end{enumerate}
		\begin{exampleblock}{Answer: D}\end{exampleblock}
		\begin{block}{Explanation}
			Training sample not representative of target population = sampling bias.
		\end{block}
	\end{frame}
	
	\begin{frame}{Q88 --- Proxy Discrimination}
		\begin{block}{Question}
			Removing race from features:
		\end{block}
		\begin{enumerate}
			\item Fully eliminates bias
			\item Can leave proxies that retain discrimination
			\item Increases bias
			\item Always legal
		\end{enumerate}
		\begin{exampleblock}{Answer: B}\end{exampleblock}
		\begin{block}{Explanation}
			Correlated features can act as proxies, preserving discrimination.
		\end{block}
	\end{frame}
	
	\begin{frame}{Q89 --- Equal Opportunity}
		\begin{block}{Question}
			Equal opportunity requires:
		\end{block}
		\begin{enumerate}
			\item Equal positive prediction rates
			\item Equal true positive rates across groups
			\item Equal precision
			\item Equal accuracy
		\end{enumerate}
		\begin{exampleblock}{Answer: B}\end{exampleblock}
		\begin{block}{Explanation}
			Equal opportunity = equal TPR across groups.
		\end{block}
	\end{frame}
	
	\begin{frame}{Q90 --- Fairness Impossibility}
		\begin{block}{Question}
			When base rates differ, which can hold simultaneously?
		\end{block}
		\begin{enumerate}
			\item All fairness metrics
			\item None of demographic parity, predictive rate parity, equal opportunity
			\item Only accuracy
			\item Only precision
		\end{enumerate}
		\begin{exampleblock}{Answer: B}\end{exampleblock}
		\begin{block}{Explanation}
			Kleinberg et al. showed these three cannot simultaneously hold when base rates differ.
		\end{block}
	\end{frame}
	
	\begin{frame}{Q91 --- Automation Bias}
		\begin{block}{Question}
			Human over-reliance on AI recommendations is called:
		\end{block}
		\begin{enumerate}
			\item Sampling bias
			\item Automation bias
			\item Historical bias
			\item Measurement bias
		\end{enumerate}
		\begin{exampleblock}{Answer: B}\end{exampleblock}
		\begin{block}{Explanation}
			Automation bias = tendency to over-trust automated systems.
		\end{block}
	\end{frame}
	
	\begin{frame}{Q92 --- Reputational Risk}
		\begin{block}{Question}
			Biggest sources of AI reputational risk:
		\end{block}
		\begin{enumerate}
			\item Cost overruns
			\item Privacy breaches, algorithmic bias, lack of explainability
			\item Slow training
			\item Too many features
		\end{enumerate}
		\begin{exampleblock}{Answer: B}\end{exampleblock}
		\begin{block}{Explanation}
			Privacy, bias, and explainability issues are the main drivers of AI reputational damage.
		\end{block}
	\end{frame}
	
	\begin{frame}{Q93 --- Misinformation}
		\begin{block}{Question}
			GenAI's biggest societal risk is:
		\end{block}
		\begin{enumerate}
			\item High compute cost
			\item Misinformation at scale and deepfakes
			\item Slow training
			\item Privacy only
		\end{enumerate}
		\begin{exampleblock}{Answer: B}\end{exampleblock}
		\begin{block}{Explanation}
			Synthetic content at scale threatens democratic institutions and public trust.
		\end{block}
	\end{frame}
	
	\begin{frame}{Q94 --- Deepfake Fraud}
		\begin{block}{Question}
			Voice deepfake CEO fraud is an example of:
		\end{block}
		\begin{enumerate}
			\item Societal risk only
			\item Impersonation risk (societal + organizational)
			\item Measurement bias
			\item Historical bias
		\end{enumerate}
		\begin{exampleblock}{Answer: B}\end{exampleblock}
		\begin{block}{Explanation}
			Deepfake fraud harms both individuals/organizations and society at large.
		\end{block}
	\end{frame}
	
	\begin{frame}{Q95 --- Hallucination}
		\begin{block}{Question}
			A chatbot confidently states wrong facts. This is:
		\end{block}
		\begin{enumerate}
			\item Bias
			\item Hallucination
			\item Overfitting
			\item Underfitting
		\end{enumerate}
		\begin{exampleblock}{Answer: B}\end{exampleblock}
		\begin{block}{Explanation}
			Hallucination = plausible but incorrect output from a probabilistic model.
		\end{block}
	\end{frame}
	
	\begin{frame}{Q96 --- Human Autonomy}
		\begin{block}{Question}
			An AI treatment recommendation conflicts with the patient's stated preferences. Which principle is most directly violated?
		\end{block}
		\begin{enumerate}
			\item Explainability
			\item Beneficence
			\item Autonomy
			\item Nonmaleficence
		\end{enumerate}
		\begin{exampleblock}{Answer: C}\end{exampleblock}
		\begin{block}{Explanation}
			Autonomy = respecting the patient's right to decide.
		\end{block}
	\end{frame}
	
	\begin{frame}{Q97 --- AI Wellbeing}
		\begin{block}{Question}
			Which risk is most associated with AI-driven mental health chatbots?
		\end{block}
		\begin{enumerate}
			\item Overfitting
			\item Lack of genuine empathy
			\item Slow response
			\item Too much accuracy
		\end{enumerate}
		\begin{exampleblock}{Answer: B}\end{exampleblock}
		\begin{block}{Explanation}
			AI can mimic empathy but lacks it. Patient outcomes drop when they know it's AI.
		\end{block}
	\end{frame}
	
	\begin{frame}{Q98 --- Job Displacement}
		\begin{block}{Question}
			Which sector is most exposed to AI automation according to recent studies?
		\end{block}
		\begin{enumerate}
			\item Manual labour only
			\item Low-skill only
			\item Both low-skill and high-skill (e.g. accounting, legal)
			\item Neither
		\end{enumerate}
		\begin{exampleblock}{Answer: C}\end{exampleblock}
		\begin{block}{Explanation}
			AI affects both low-skill and high-skill cognitive tasks.
		\end{block}
	\end{frame}
	
	\begin{frame}{Q99 --- Feedback Loop}
		\begin{block}{Question}
			A hiring model trained on historical data perpetuates past discrimination. This is an example of:
		\end{block}
		\begin{enumerate}
			\item Measurement bias
			\item Feedback loop
			\item Sampling noise
			\item Deployment error
		\end{enumerate}
		\begin{exampleblock}{Answer: B}\end{exampleblock}
		\begin{block}{Explanation}
			Feedback loops amplify historical patterns, reinforcing discrimination.
		\end{block}
	\end{frame}
	
	\begin{frame}{Q100 --- Group Fairness}
		\begin{block}{Question}
			Which of the following is NOT a measure of group fairness?
		\end{block}
		\begin{enumerate}
			\item Demographic parity
			\item Predictive rate parity
			\item Equal opportunity
			\item Model accuracy
		\end{enumerate}
		\begin{exampleblock}{Answer: D}\end{exampleblock}
		\begin{block}{Explanation}
			Accuracy is a performance metric, not a fairness measure.
		\end{block}
	\end{frame}
	
	% =================================================================
	\section{Module 4: Responsible \& Ethical AI}
	% =================================================================
	
	\begin{frame}{Q101 --- Shapley Values}
		\begin{block}{Question}
			Which technique best explains the marginal effect of each feature?
		\end{block}
		\begin{enumerate}
			\item Feature importance
			\item Shapley values
			\item Grid search
			\item Feature scaling
		\end{enumerate}
		\begin{exampleblock}{Answer: B}\end{exampleblock}
		\begin{block}{Explanation}
			Shapley values calculate marginal contribution of each feature across all subsets.
		\end{block}
	\end{frame}
	
	\begin{frame}{Q102 --- Loan Explainability}
		\begin{block}{Question}
			Customers want to understand why their loans were rejected. Which is the source of reputational risk?
		\end{block}
		\begin{enumerate}
			\item Privacy breaches
			\item Algorithmic bias
			\item Lack of explainability
			\item Lack of transparency
		\end{enumerate}
		\begin{exampleblock}{Answer: C}\end{exampleblock}
		\begin{block}{Explanation}
			Provide reasons for rejections; otherwise customers lose trust.
		\end{block}
	\end{frame}
	
	\begin{frame}{Q103 --- Practical Ethics}
		\begin{block}{Question}
			What is the advantage of a practical ethics framework?
		\end{block}
		\begin{enumerate}
			\item Provides clear guidelines for ethical dilemmas
			\item Eliminates conflicts of interest
			\item Guarantees legal compliance
			\item Focuses only on profit
		\end{enumerate}
		\begin{exampleblock}{Answer: A}\end{exampleblock}
		\begin{block}{Explanation}
			Provides structured guidelines. It cannot eliminate conflicts or guarantee compliance.
		\end{block}
	\end{frame}
	
	\begin{frame}{Q104 --- Ethics Framework Benefit}
		\begin{block}{Question}
			What is the most likely benefit of implementing an AI ethics framework?
		\end{block}
		\begin{enumerate}
			\item Increases model accuracy
			\item Strengthens firewalls
			\item Increases productivity
			\item Reduces risk of regulatory non-compliance
		\end{enumerate}
		\begin{exampleblock}{Answer: D}\end{exampleblock}
		\begin{block}{Explanation}
			Ethics frameworks align AI with legal and ethical standards, reducing regulatory risk.
		\end{block}
	\end{frame}
	
	\begin{frame}{Q105 --- Utilitarianism}
		\begin{block}{Question}
			Which statement is correct regarding utilitarianism?
		\end{block}
		\begin{enumerate}
			\item Can lead to violations of individual rights
			\item Focuses on virtuous character
			\item Relies heavily on laws
			\item Values intrinsic morality of actions
		\end{enumerate}
		\begin{exampleblock}{Answer: A}\end{exampleblock}
		\begin{block}{Explanation}
			Utilitarianism maximises overall good, sometimes at the cost of individual rights.
		\end{block}
	\end{frame}
	
	\begin{frame}{Q106 --- Deontology}
		\begin{block}{Question}
			Which best describes deontology?
		\end{block}
		\begin{enumerate}
			\item Judges morality by consequences
			\item Judges morality by duties and rules
			\item Uses a single quantifiable metric
			\item Emphasises virtuous character
		\end{enumerate}
		\begin{exampleblock}{Answer: B}\end{exampleblock}
		\begin{block}{Explanation}
			Deontology judges by adherence to duties and rules, not consequences.
		\end{block}
	\end{frame}
	
	\begin{frame}{Q107 --- Virtue Ethics}
		\begin{block}{Question}
			Which is a defining feature of virtue ethics?
		\end{block}
		\begin{enumerate}
			\item Rules and duties
			\item Maximising utility
			\item Character traits and living a good life
			\item Legal compliance
		\end{enumerate}
		\begin{exampleblock}{Answer: C}\end{exampleblock}
		\begin{block}{Explanation}
			Virtue ethics focuses on character and virtues.
		\end{block}
	\end{frame}
	
	\begin{frame}{Q108 --- Justice}
		\begin{block}{Question}
			Which statement is correct regarding the principle of justice in AI ethics?
		\end{block}
		\begin{enumerate}
			\item Conflicts can exist between different concepts of distributive justice
			\item Impartiality is not required for procedural justice
			\item Autonomy is secondary to social justice
			\item Meritocracy is not a distributive concept
		\end{enumerate}
		\begin{exampleblock}{Answer: A}\end{exampleblock}
		\begin{block}{Explanation}
			Egalitarianism, utilitarianism, and meritocracy can conflict.
		\end{block}
	\end{frame}
	
	\begin{frame}{Q109 --- Autonomy}
		\begin{block}{Question}
			An AI recommends treatment conflicting with the patient's preferences. Which principle is violated?
		\end{block}
		\begin{enumerate}
			\item Explainability
			\item Beneficence
			\item Autonomy
			\item Nonmaleficence
		\end{enumerate}
		\begin{exampleblock}{Answer: C}\end{exampleblock}
		\begin{block}{Explanation}
			Autonomy = respecting the patient's own decisions.
		\end{block}
	\end{frame}
	
	\begin{frame}{Q110 --- Nonmaleficence}
		\begin{block}{Question}
			Which principle is best described as ``do no harm''?
		\end{block}
		\begin{enumerate}
			\item Beneficence
			\item Nonmaleficence
			\item Justice
			\item Autonomy
		\end{enumerate}
		\begin{exampleblock}{Answer: B}\end{exampleblock}
		\begin{block}{Explanation}
			Nonmaleficence = avoid harm, including unjustified risk.
		\end{block}
	\end{frame}
	
	\begin{frame}{Q111 --- Beneficence}
		\begin{block}{Question}
			Which principle commands actively promoting the welfare of others?
		\end{block}
		\begin{enumerate}
			\item Nonmaleficence
			\item Beneficence
			\item Justice
			\item Autonomy
		\end{enumerate}
		\begin{exampleblock}{Answer: B}\end{exampleblock}
		\begin{block}{Explanation}
			Beneficence = positive action to help others.
		\end{block}
	\end{frame}
	
	\begin{frame}{Q112 --- Privacy}
		\begin{block}{Question}
			What is the most effective way to minimize privacy risk?
		\end{block}
		\begin{enumerate}
			\item Collect all data regardless of sensitivity
			\item Data minimization and giving control to data subjects
			\item Share data with third parties freely
			\item Store unencrypted backups
		\end{enumerate}
		\begin{exampleblock}{Answer: B}\end{exampleblock}
		\begin{block}{Explanation}
			Data minimization and user control reduce exposure.
		\end{block}
	\end{frame}
	
	\begin{frame}{Q113 --- Right to be Forgotten}
		\begin{block}{Question}
			The GDPR right to be forgotten allows individuals to:
		\end{block}
		\begin{enumerate}
			\item Access their data
			\item Request deletion of personal data
			\item Sell their data
			\item Share data without consent
		\end{enumerate}
		\begin{exampleblock}{Answer: B}\end{exampleblock}
		\begin{block}{Explanation}
			Right to erasure = deletion of personal data under certain conditions.
		\end{block}
	\end{frame}
	
	\begin{frame}{Q114 --- Contextual Integrity}
		\begin{block}{Question}
			Data given to a bank for a transaction should:
		\end{block}
		\begin{enumerate}
			\item Be sold freely
			\item Only be used according to contextual norms
			\item Be publicly shared
			\item Be kept forever
		\end{enumerate}
		\begin{exampleblock}{Answer: B}\end{exampleblock}
		\begin{block}{Explanation}
			Contextual integrity: data use must respect expectations of the context.
		\end{block}
	\end{frame}
	
	\begin{frame}{Q115 --- EU AI Act}
		\begin{block}{Question}
			The EU AI Act prohibits:
		\end{block}
		\begin{enumerate}
			\item All AI systems
			\item Social scoring and manipulative AI
			\item Only chatbots
			\item Medical AI
		\end{enumerate}
		\begin{exampleblock}{Answer: B}\end{exampleblock}
		\begin{block}{Explanation}
			Prohibited: social scoring, manipulative AI, exploitation of vulnerabilities.
		\end{block}
	\end{frame}
	
	\begin{frame}{Q116 --- GDPR}
		\begin{block}{Question}
			What is the key feature of GDPR?
		\end{block}
		\begin{enumerate}
			\item Applies only to EU companies
			\item Extraterritorial jurisdiction for services to EU residents
			\item No fines
			\item Voluntary compliance
		\end{enumerate}
		\begin{exampleblock}{Answer: B}\end{exampleblock}
		\begin{block}{Explanation}
			GDPR applies globally if serving EU data subjects.
		\end{block}
	\end{frame}
	
	\begin{frame}{Q117 --- NIST AI RMF}
		\begin{block}{Question}
			The NIST AI RMF is structured around four functions:
		\end{block}
		\begin{enumerate}
			\item Govern, map, measure, manage
			\item Train, test, deploy, monitor
			\item Audit, report, remediate, retire
			\item Design, build, test, launch
		\end{enumerate}
		\begin{exampleblock}{Answer: A}\end{exampleblock}
		\begin{block}{Explanation}
			Govern, map, measure, manage.
		\end{block}
	\end{frame}
	
	\begin{frame}{Q118 --- Utilitarianism Founders}
		\begin{block}{Question}
			Who are the classical proponents of utilitarianism?
		\end{block}
		\begin{enumerate}
			\item Kant and Aristotle
			\item Bentham and Mill
			\item Rawls and Nozick
			\item Plato and Socrates
		\end{enumerate}
		\begin{exampleblock}{Answer: B}\end{exampleblock}
		\begin{block}{Explanation}
			Jeremy Bentham and John Stuart Mill developed utilitarianism.
		\end{block}
	\end{frame}
	
	\begin{frame}{Q119 --- Kant's Categorical Imperative}
		\begin{block}{Question}
			What does Kant's categorical imperative require?
		\end{block}
		\begin{enumerate}
			\item Maximise utility
			\item Act on maxims that could be universal laws
			\item Maximise virtue
			\item Maximise profit
		\end{enumerate}
		\begin{exampleblock}{Answer: B}\end{exampleblock}
		\begin{block}{Explanation}
			Universalizability: act only on maxims you would will as universal law.
		\end{block}
	\end{frame}
	
	\begin{frame}{Q120 --- Medical Ethics Origins}
		\begin{block}{Question}
			Which two principles of AI ethics derive from medical ethics?
		\end{block}
		\begin{enumerate}
			\item Marketing and profit
			\item Nonmaleficence and beneficence
			\item Privacy and censorship
			\item Efficiency and speed
		\end{enumerate}
		\begin{exampleblock}{Answer: B}\end{exampleblock}
		\begin{block}{Explanation}
			Nonmaleficence and beneficence are central to both medical and AI ethics.
		\end{block}
	\end{frame}
	
	\begin{frame}{Q121 --- Privacy as Wrong}
		\begin{block}{Question}
			Violating privacy is morally wrong even without tangible harm because:
		\end{block}
		\begin{enumerate}
			\item It costs money
			\item It violates rights and treats people as means
			\item It slows training
			\item It confuses regulators
		\end{enumerate}
		\begin{exampleblock}{Answer: B}\end{exampleblock}
		\begin{block}{Explanation}
			Privacy violations wrong individuals by treating them instrumentally.
		\end{block}
	\end{frame}
	
	\begin{frame}{Q122 --- Data Minimization}
		\begin{block}{Question}
			Which practice is most effective for privacy protection?
		\end{block}
		\begin{enumerate}
			\item Collect all possible data
			\item Data minimization
			\item Share data with vendors
			\item Store in plaintext
		\end{enumerate}
		\begin{exampleblock}{Answer: B}\end{exampleblock}
		\begin{block}{Explanation}
			Collecting only what is needed reduces privacy risk.
		\end{block}
	\end{frame}
	
	\begin{frame}{Q123 --- Synthetic Data}
		\begin{block}{Question}
			What is an advantage of synthetic data?
		\end{block}
		\begin{enumerate}
			\item Always more accurate
			\item Reduced privacy risk
			\item Always faster to generate
			\item No cost
		\end{enumerate}
		\begin{exampleblock}{Answer: B}\end{exampleblock}
		\begin{block}{Explanation}
			Synthetic data doesn't come from real subjects, reducing privacy risk.
		\end{block}
	\end{frame}
	
	\begin{frame}{Q124 --- Algorithmic Auditing}
		\begin{block}{Question}
			What is the best way to detect algorithmic bias?
		\end{block}
		\begin{enumerate}
			\item Ignore it
			\item Regular, independent audits
			\item Only use open-source models
			\item Remove all sensitive features
		\end{enumerate}
		\begin{exampleblock}{Answer: B}\end{exampleblock}
		\begin{block}{Explanation}
			Independent algorithmic audits are the most reliable method.
		\end{block}
	\end{frame}
	
	\begin{frame}{Q125 --- Fairness Automation}
		\begin{block}{Question}
			Why can't fairness be fully automated?
		\end{block}
		\begin{enumerate}
			\item It's too expensive
			\item Unequal base rates make simultaneous parity mathematically impossible
			\item Not enough data
			\item GDPR forbids it
		\end{enumerate}
		\begin{exampleblock}{Answer: B}\end{exampleblock}
		\begin{block}{Explanation}
			Different fairness metrics cannot all hold when base rates differ.
		\end{block}
	\end{frame}
	
	\begin{frame}{Q126 --- Explainability Levels}
		\begin{block}{Question}
			Local explainability differs from global explainability by:
		\end{block}
		\begin{enumerate}
			\item Model type
			\item Scope: single prediction vs whole-model behaviour
			\item Speed
			\item Data size
		\end{enumerate}
		\begin{exampleblock}{Answer: B}\end{exampleblock}
		\begin{block}{Explanation}
			Local = specific decision; global = overall behaviour.
		\end{block}
	\end{frame}
	
	\begin{frame}{Q127 --- Counterfactual Explanation}
		\begin{block}{Question}
			``If income were \$10k higher, the loan would have been approved.'' This is a:
		\end{block}
		\begin{enumerate}
			\item SHAP value
			\item Counterfactual explanation
			\item LIME
			\item Feature importance
		\end{enumerate}
		\begin{exampleblock}{Answer: B}\end{exampleblock}
		\begin{block}{Explanation}
			Counterfactuals describe alternative inputs that change the outcome.
		\end{block}
	\end{frame}
	
	\begin{frame}{Q128 --- LIME}
		\begin{block}{Question}
			What does LIME provide?
		\end{block}
		\begin{enumerate}
			\item Global feature ranking
			\item Local surrogate explanations
			\item Fairness guarantees
			\item Privacy preservation
		\end{enumerate}
		\begin{exampleblock}{Answer: B}\end{exampleblock}
		\begin{block}{Explanation}
			LIME approximates the model locally with an interpretable surrogate.
		\end{block}
	\end{frame}
	
	\begin{frame}{Q129 --- SHAP}
		\begin{block}{Question}
			SHAP values are grounded in:
		\end{block}
		\begin{enumerate}
			\item Bayesian inference
			\item Marginal contribution theory from cooperative game theory
			\item Gradient descent
			\item Decision tree splitting
		\end{enumerate}
		\begin{exampleblock}{Answer: B}\end{exampleblock}
		\begin{block}{Explanation}
			SHAP = SHapley Additive exPlanations, based on Shapley values.
		\end{block}
	\end{frame}
	
	\begin{frame}{Q130 --- Accuracy/Interpretability}
		\begin{block}{Question}
			Which statement about the accuracy-interpretability trade-off is correct?
		\end{block}
		\begin{enumerate}
			\item Always choose accuracy
			\item Complex models are often more accurate but less interpretable
			\item Simple models always win
			\item Interpretability is irrelevant
		\end{enumerate}
		\begin{exampleblock}{Answer: B}\end{exampleblock}
		\begin{block}{Explanation}
			Complex models (neural nets) often outperform simple models but are harder to explain.
		\end{block}
	\end{frame}
	
	\begin{frame}{Q131 --- EU AI Act Fines}
		\begin{block}{Question}
			What is the maximum fine under the EU AI Act?
		\end{block}
		\begin{enumerate}
			\item €1 million
			\item €10 million
			\item €35 million or 7\% of global turnover
			\item No fines
		\end{enumerate}
		\begin{exampleblock}{Answer: C}\end{exampleblock}
		\begin{block}{Explanation}
			Up to €35 million or 7\% of global turnover for the most serious violations.
		\end{block}
	\end{frame}
	
	\begin{frame}{Q132 --- CCPA}
		\begin{block}{Question}
			The CCPA gives California consumers the right to:
		\end{block}
		\begin{enumerate}
			\item Know what data is collected
			\item Opt out of the sale of personal information
			\item Delete personal information
			\item All of the above
		\end{enumerate}
		\begin{exampleblock}{Answer: D}\end{exampleblock}
		\begin{block}{Explanation}
			CCPA provides knowledge, opt-out, and deletion rights.
		\end{block}
	\end{frame}
	
	\begin{frame}{Q133 --- GDPR Data Breach}
		\begin{block}{Question}
			Under GDPR, organizations must notify the supervisory authority of a data breach within:
		\end{block}
		\begin{enumerate}
			\item 24 hours
			\item 48 hours
			\item 72 hours
			\item 7 days
		\end{enumerate}
		\begin{exampleblock}{Answer: C}\end{exampleblock}
		\begin{block}{Explanation}
			72 hours is the GDPR requirement.
		\end{block}
	\end{frame}
	
	\begin{frame}{Q134 --- DSA}
		\begin{block}{Question}
			The EU Digital Services Act requires platforms to:
		\end{block}
		\begin{enumerate}
			\item Prevent dissemination of illegal content
			\item Be transparent about content moderation
			\item Address disinformation
			\item All of the above
		\end{enumerate}
		\begin{exampleblock}{Answer: D}\end{exampleblock}
		\begin{block}{Explanation}
			DSA covers illegal content, transparency, disinformation, and algorithmic bias.
		\end{block}
	\end{frame}
	
	\begin{frame}{Q135 --- Utilitarian Critique}
		\begin{block}{Question}
			A key criticism of utilitarianism is that it:
		\end{block}
		\begin{enumerate}
			\item Ignores consequences
			\item Can justify violating individual rights for greater good
			\item Emphasises character
			\item Relies only on rules
		\end{enumerate}
		\begin{exampleblock}{Answer: B}\end{exampleblock}
		\begin{block}{Explanation}
			Maximising total utility can justify serious rights violations.
		\end{block}
	\end{frame}
	
	\begin{frame}{Q136 --- Virtue Ethics Critics}
		\begin{block}{Question}
			Critics of virtue ethics argue that it:
		\end{block}
		\begin{enumerate}
			\item Lacks clear guidance for moral decisions
			\item Ignores character
			\item Is too rule-based
			\item Only considers consequences
		\end{enumerate}
		\begin{exampleblock}{Answer: A}\end{exampleblock}
		\begin{block}{Explanation}
			Critics say virtue ethics lacks clear decision procedures.
		\end{block}
	\end{frame}
	
	\begin{frame}{Q137 --- Medical Ethics Timeline}
		\begin{block}{Question}
			Which document was foundational for medical ethics in 1947?
		\end{block}
		\begin{enumerate}
			\item Nuremberg Code
			\item HIPAA
			\item GDPR
			\item Belmont Report
		\end{enumerate}
		\begin{exampleblock}{Answer: A}\end{exampleblock}
		\begin{block}{Explanation}
			The Nuremberg Code (1947) established ethical research principles.
		\end{block}
	\end{frame}
	
	\begin{frame}{Q138 --- AI Ethics Principles}
		\begin{block}{Question}
			Which is NOT one of the five commonly cited AI ethics principles?
		\end{block}
		\begin{enumerate}
			\item Nonmaleficence
			\item Beneficence
			\item Profit maximization
			\item Autonomy
		\end{enumerate}
		\begin{exampleblock}{Answer: C}\end{exampleblock}
		\begin{block}{Explanation}
			The five: nonmaleficence, beneficence, justice, autonomy, explainability.
		\end{block}
	\end{frame}
	
	\begin{frame}{Q139 --- Beneficence Limit}
		\begin{block}{Question}
			Which limits the duty of beneficence?
			\begin{enumerate}
				\item Unlimited resources
				\item Scarce resources and competing obligations
				\item No limits
				\item Only legal constraints
			\end{enumerate}
		\end{block}
		\begin{exampleblock}{Answer: B}\end{exampleblock}
		\begin{block}{Explanation}
			Beneficence is constrained by reasonableness and scarce resources.
		\end{block}
	\end{frame}
	
	\begin{frame}{Q140 --- Autonomy Surrogate}
		\begin{block}{Question}
			When a person cannot make autonomous decisions, who should decide?
		\end{block}
		\begin{enumerate}
			\item The AI system
			\item A surrogate decision-maker acting in the person's best interest
			\item The government
			\item No one
		\end{enumerate}
		\begin{exampleblock}{Answer: B}\end{exampleblock}
		\begin{block}{Explanation}
			Surrogates (family, guardians) decide for non-autonomous individuals.
		\end{block}
	\end{frame}
	
	% =================================================================
	\section{Module 5: Data and AI Model Governance}
	% =================================================================
	
	\begin{frame}{Q141 --- Model Landscape}
		\begin{block}{Question}
			What is a crucial consideration when evaluating the model landscape?
		\end{block}
		\begin{enumerate}
			\item Verify each model operates independently
			\item Ensure all models use the same input data
			\item Assess whether models can be consolidated
			\item Ensure interdependencies between models are identified
		\end{enumerate}
		\begin{exampleblock}{Answer: D}\end{exampleblock}
		\begin{block}{Explanation}
			The landscape shows interactions; identifying dependencies is critical.
		\end{block}
	\end{frame}
	
	\begin{frame}{Q142 --- Governance Roles}
		\begin{block}{Question}
			Which correctly describes model governance roles?
		\end{block}
		\begin{enumerate}
			\item Board maintains inventory
			\item Model risk function sets appetite
			\item Auditors are independent checks on validation
			\item Ownership never shared
		\end{enumerate}
		\begin{exampleblock}{Answer: C}\end{exampleblock}
		\begin{block}{Explanation}
			Auditors (internal/external) provide independent oversight of validation.
		\end{block}
	\end{frame}
	
	\begin{frame}{Q143 --- IT Integration}
		\begin{block}{Question}
			Which team ensures model integration with IT systems does not disrupt operations?
		\end{block}
		\begin{enumerate}
			\item Core Implementation Team
			\item System Implementation Team
			\item Model Design Team
			\item Data Management Team
		\end{enumerate}
		\begin{exampleblock}{Answer: B}\end{exampleblock}
		\begin{block}{Explanation}
			System Implementation Team handles IT integration and production deployment.
		\end{block}
	\end{frame}
	
	\begin{frame}{Q144 --- Validation Next Step}
		\begin{block}{Question}
			A new model outperforms the traditional one. What is the appropriate next step?
		\end{block}
		\begin{enumerate}
			\item Deploy to production
			\item Inform IT
			\item Retire the existing model
			\item Request initial validation
		\end{enumerate}
		\begin{exampleblock}{Answer: D}\end{exampleblock}
		\begin{block}{Explanation}
			Validation is required before deployment.
		\end{block}
	\end{frame}
	
	\begin{frame}{Q145 --- Alternative Data}
		\begin{block}{Question}
			A credit model uses school transcripts, social media, landlord disputes. What is an appropriate validation finding?
		\end{block}
		\begin{enumerate}
			\item No concerns
			\item Raise concerns about sensitive alternative data
			\item Approve if backtest improves
			\item Add more alternative data
		\end{enumerate}
		\begin{exampleblock}{Answer: B}\end{exampleblock}
		\begin{block}{Explanation}
			Sensitive alternative data requires heightened scrutiny and consent.
		\end{block}
	\end{frame}
	
	\begin{frame}{Q146 --- Monitoring Frequency}
		\begin{block}{Question}
			What is an advantage of increasing model monitoring frequency?
		\end{block}
		\begin{enumerate}
			\item Reduces need for retraining
			\item Allows detection of performance degradation or drift
			\item Minimizes IT involvement
			\item Improves interpretability
		\end{enumerate}
		\begin{exampleblock}{Answer: B}\end{exampleblock}
		\begin{block}{Explanation}
			Frequent monitoring detects drift and enables timely intervention.
		\end{block}
	\end{frame}
	
	\begin{frame}{Q147 --- Model Inventory}
		\begin{block}{Question}
			What is the purpose of a model inventory?
		\end{block}
		\begin{enumerate}
			\item Store model code
			\item List all models and track their risk
			\item Train models
			\item Visualize performance
		\end{enumerate}
		\begin{exampleblock}{Answer: B}\end{exampleblock}
		\begin{block}{Explanation}
			Inventory = registry of all models with purpose, risk, owner, validation status.
		\end{block}
	\end{frame}
	
	\begin{frame}{Q148 --- Documentation}
		\begin{block}{Question}
			What is the primary purpose of model documentation?
		\end{block}
		\begin{enumerate}
			\item Increase accuracy
			\item Record design, assumptions, and limitations
			\item Speed up training
			\item Reduce features
		\end{enumerate}
		\begin{exampleblock}{Answer: B}\end{exampleblock}
		\begin{block}{Explanation}
			Documentation supports validation, audit, knowledge transfer, compliance.
		\end{block}
	\end{frame}
	
	\begin{frame}{Q149 --- Use Test}
		\begin{block}{Question}
			What is the purpose of the use test in model validation?
		\end{block}
		\begin{enumerate}
			\item Test production use
			\item Assess appropriateness for intended use
			\item Test speed
			\item Test accuracy
		\end{enumerate}
		\begin{exampleblock}{Answer: B}\end{exampleblock}
		\begin{block}{Explanation}
			Use test assesses fit for intended purpose and user understanding.
		\end{block}
	\end{frame}
	
	\begin{frame}{Q150 --- White vs Black Box}
		\begin{block}{Question}
			What is the difference between white-box and black-box testing?
		\end{block}
		\begin{enumerate}
			\item White-box tests internal logic; black-box tests inputs/outputs
			\item White-box tests inputs/outputs; black-box tests internal logic
			\item They are the same
			\item White-box for AI; black-box for statistics
		\end{enumerate}
		\begin{exampleblock}{Answer: A}\end{exampleblock}
		\begin{block}{Explanation}
			White-box: internal code/logic. Black-box: behaviour only.
		\end{block}
	\end{frame}
	
	\begin{frame}{Q151 --- Data Provenance}
		\begin{block}{Question}
			What is the purpose of data provenance?
		\end{block}
		\begin{enumerate}
			\item Increase data volume
			\item Verify the legal right to use data
			\item Speed up training
			\item Reduce data quality
		\end{enumerate}
		\begin{exampleblock}{Answer: B}\end{exampleblock}
		\begin{block}{Explanation}
			Provenance = origin and legal right to use data.
		\end{block}
	\end{frame}
	
	\begin{frame}{Q152 --- Data Classification}
		\begin{block}{Question}
			What is the primary purpose of data classification?
		\end{block}
		\begin{enumerate}
			\item Increase volume
			\item Categorize by sensitivity and structure
			\item Speed up training
			\item Reduce features
		\end{enumerate}
		\begin{exampleblock}{Answer: B}\end{exampleblock}
		\begin{block}{Explanation}
			Classification determines handling, access controls, protection.
		\end{block}
	\end{frame}
	
	\begin{frame}{Q153 --- Metadata}
		\begin{block}{Question}
			Metadata is best described as:
		\end{block}
		\begin{enumerate}
			\item Raw data
			\item Data about data
			\item Encrypted data
			\item Synthetic data
		\end{enumerate}
		\begin{exampleblock}{Answer: B}\end{exampleblock}
		\begin{block}{Explanation}
			Metadata describes origin, structure, definitions, relationships.
		\end{block}
	\end{frame}
	
	\begin{frame}{Q154 --- Access Controls}
		\begin{block}{Question}
			What is the purpose of data access controls?
		\end{block}
		\begin{enumerate}
			\item Increase volume
			\item Ensure only authorized users can access data
			\item Speed up training
			\item Reduce quality
		\end{enumerate}
		\begin{exampleblock}{Answer: B}\end{exampleblock}
		\begin{block}{Explanation}
			Access controls protect privacy and security.
		\end{block}
	\end{frame}
	
	\begin{frame}{Q155 --- Data Governance}
		\begin{block}{Question}
			What is the primary purpose of a data governance framework?
		\end{block}
		\begin{enumerate}
			\item Increase data volume
			\item Ensure data quality, security, and compliance
			\item Speed up training
			\item Reduce features
		\end{enumerate}
		\begin{exampleblock}{Answer: B}\end{exampleblock}
		\begin{block}{Explanation}
			Data governance ensures quality, security, compliance, proper use.
		\end{block}
	\end{frame}
	
	\begin{frame}{Q156 --- Model Risk Management}
		\begin{block}{Question}
			What is the purpose of model risk management?
		\end{block}
		\begin{enumerate}
			\item Increase accuracy
			\item Identify, assess, and mitigate risks from models
			\item Speed up training
			\item Reduce features
		\end{enumerate}
		\begin{exampleblock}{Answer: B}\end{exampleblock}
		\begin{block}{Explanation}
			MRM identifies, assesses, mitigates model risks.
		\end{block}
	\end{frame}
	
	\begin{frame}{Q157 --- Model Risk Function}
		\begin{block}{Question}
			What is the role of the model risk function?
		\end{block}
		\begin{enumerate}
			\item Build models
			\item Maintain inventory and develop policies
			\item Use models
			\item Reduce features
		\end{enumerate}
		\begin{exampleblock}{Answer: B}\end{exampleblock}
		\begin{block}{Explanation}
			Model risk function = inventory + policies + oversight.
		\end{block}
	\end{frame}
	
	\begin{frame}{Q158 --- Periodic Validation}
		\begin{block}{Question}
			What is the purpose of periodic model validation?
		\end{block}
		\begin{enumerate}
			\item Increase accuracy
			\item Confirm the model remains fit for purpose over time
			\item Speed up training
			\item Reduce features
		\end{enumerate}
		\begin{exampleblock}{Answer: B}\end{exampleblock}
		\begin{block}{Explanation}
			Periodic validation confirms assumptions, data, performance remain valid.
		\end{block}
	\end{frame}
	
	\begin{frame}{Q159 --- Implementation}
		\begin{block}{Question}
			What is the purpose of model implementation tasks?
		\end{block}
		\begin{enumerate}
			\item Design the model
			\item Deploy to production and integrate with IT
			\item Validate the model
			\item Reduce features
		\end{enumerate}
		\begin{exampleblock}{Answer: B}\end{exampleblock}
		\begin{block}{Explanation}
			Implementation = deployment and IT integration.
		\end{block}
	\end{frame}
	
	\begin{frame}{Q160 --- Adaptation}
		\begin{block}{Question}
			What is the purpose of model adaptation tasks?
		\end{block}
		\begin{enumerate}
			\item Design the model
			\item Update the model for changing conditions
			\item Validate the model
			\item Reduce features
		\end{enumerate}
		\begin{exampleblock}{Answer: B}\end{exampleblock}
		\begin{block}{Explanation}
			Adaptation = retraining, recalibration, updating assumptions.
		\end{block}
	\end{frame}
	
	\begin{frame}{Q161 --- Risk Appetite}
		\begin{block}{Question}
			Who sets the model risk appetite?
		\end{block}
		\begin{enumerate}
			\item Model developers
			\item Board of directors
			\item Auditors
			\item Users
		\end{enumerate}
		\begin{exampleblock}{Answer: B}\end{exampleblock}
		\begin{block}{Explanation}
			Board sets risk appetite; model risk function maintains policies.
		\end{block}
	\end{frame}
	
	\begin{frame}{Q162 --- Policies}
		\begin{block}{Question}
			What is the purpose of model governance policies?
		\end{block}
		\begin{enumerate}
			\item Increase accuracy
			\item Establish rules for development, validation, use
			\item Speed up training
			\item Reduce features
		\end{enumerate}
		\begin{exampleblock}{Answer: B}\end{exampleblock}
		\begin{block}{Explanation}
			Policies provide consistency, accountability, compliance.
		\end{block}
	\end{frame}
	
	\begin{frame}{Q163 --- Documentation Standards}
		\begin{block}{Question}
			What is the purpose of documentation standards?
		\end{block}
		\begin{enumerate}
			\item Increase accuracy
			\item Ensure consistent and comprehensive documentation
			\item Speed up training
			\item Reduce features
		\end{enumerate}
		\begin{exampleblock}{Answer: B}\end{exampleblock}
		\begin{block}{Explanation}
			Standards ensure consistency and audit readiness.
		\end{block}
	\end{frame}
	
	\begin{frame}{Q164 --- Inventory vs Landscape}
		\begin{block}{Question}
			What is the difference between inventory and landscape?
		\end{block}
		\begin{enumerate}
			\item Inventory is a list; landscape shows interactions
			\item They are the same
			\item Landscape is a list; inventory shows interactions
			\item Only inventory matters
		\end{enumerate}
		\begin{exampleblock}{Answer: A}\end{exampleblock}
		\begin{block}{Explanation}
			Inventory = registry. Landscape = map of interactions.
		\end{block}
	\end{frame}
	
	\begin{frame}{Q165 --- Reporting}
		\begin{block}{Question}
			What is the purpose of model risk reporting?
		\end{block}
		\begin{enumerate}
			\item Increase accuracy
			\item Inform senior management and board about model risk
			\item Speed up training
			\item Reduce features
		\end{enumerate}
		\begin{exampleblock}{Answer: B}\end{exampleblock}
		\begin{block}{Explanation}
			Reporting communicates risk to governance bodies.
		\end{block}
	\end{frame}
	
	\begin{frame}{Q166 --- Change Management}
		\begin{block}{Question}
			What is the purpose of model change management?
		\end{block}
		\begin{enumerate}
			\item Increase accuracy
			\item Control changes with review and approval
			\item Speed up training
			\item Reduce features
		\end{enumerate}
		\begin{exampleblock}{Answer: B}\end{exampleblock}
		\begin{block}{Explanation}
			Change management ensures controlled, documented changes.
		\end{block}
	\end{frame}
	
	\begin{frame}{Q167 --- Model Review}
		\begin{block}{Question}
			What is the purpose of model review?
		\end{block}
		\begin{enumerate}
			\item Increase accuracy
			\item Periodically assess model performance and fitness
			\item Speed up training
			\item Reduce features
		\end{enumerate}
		\begin{exampleblock}{Answer: B}\end{exampleblock}
		\begin{block}{Explanation}
			Review evaluates continued fitness and recommends action.
		\end{block}
	\end{frame}
	
	\begin{frame}{Q168 --- Retirement}
		\begin{block}{Question}
			What is the purpose of model retirement?
		\end{block}
		\begin{enumerate}
			\item Increase accuracy
			\item Formally decommission a model no longer needed
			\item Speed up training
			\item Reduce features
		\end{enumerate}
		\begin{exampleblock}{Answer: B}\end{exampleblock}
		\begin{block}{Explanation}
			Retirement = formal decommissioning with documentation.
		\end{block}
	\end{frame}
	
	\begin{frame}{Q169 --- Risk Culture}
		\begin{block}{Question}
			What is the purpose of model risk culture?
		\end{block}
		\begin{enumerate}
			\item Increase accuracy
			\item Promote awareness and accountability across the organization
			\item Speed up training
			\item Reduce features
		\end{enumerate}
		\begin{exampleblock}{Answer: B}\end{exampleblock}
		\begin{block}{Explanation}
			Culture promotes awareness and accountability at all levels.
		\end{block}
	\end{frame}
	
	\begin{frame}{Q170 --- Training}
		\begin{block}{Question}
			What is the purpose of model risk training?
		\end{block}
		\begin{enumerate}
			\item Increase accuracy
			\item Ensure staff understand model risk and their responsibilities
			\item Speed up training
			\item Reduce features
		\end{enumerate}
		\begin{exampleblock}{Answer: B}\end{exampleblock}
		\begin{block}{Explanation}
			Training builds awareness and competence.
		\end{block}
	\end{frame}
	
	\begin{frame}{Q171 --- Three Lines of Defense}
		\begin{block}{Question}
			Which correctly describes the Three Lines of Defense?
		\end{block}
		\begin{enumerate}
			\item First line: business; Second: independent risk; Third: internal audit
			\item First: audit; Second: business; Third: risk
			\item All three audit the same thing
			\item Only two lines exist
		\end{enumerate}
		\begin{exampleblock}{Answer: A}\end{exampleblock}
		\begin{block}{Explanation}
			First line owns the model; second line validates; third line audits.
		\end{block}
	\end{frame}
	
	\begin{frame}{Q172 --- Model Owner}
		\begin{block}{Question}
			Who is typically assigned model ownership?
		\end{block}
		\begin{enumerate}
			\item External vendors
			\item Users in the business unit
			\item Auditors
			\item Regulators
		\end{enumerate}
		\begin{exampleblock}{Answer: B}\end{exampleblock}
		\begin{block}{Explanation}
			Model owners are usually in the business unit, ensuring accountability.
		\end{block}
	\end{frame}
	
	\begin{frame}{Q173 --- Model Validators}
		\begin{block}{Question}
			Model validators must be:
		\end{block}
		\begin{enumerate}
			\item The same team as developers
			\item Independent from model development
			\item External only
			\item Junior staff
		\end{enumerate}
		\begin{exampleblock}{Answer: B}\end{exampleblock}
		\begin{block}{Explanation}
			Independence is essential for objective validation.
		\end{block}
	\end{frame}
	
	\begin{frame}{Q174 --- Model Developers}
		\begin{block}{Question}
			What role do model developers play in validation?
		\end{block}
		\begin{enumerate}
			\item They approve the model
			\item They provide documentation and respond to findings
			\item They perform independent validation
			\item They have no role
		\end{enumerate}
		\begin{exampleblock}{Answer: B}\end{exampleblock}
		\begin{block}{Explanation}
			Developers produce documentation and remediate issues.
		\end{block}
	\end{frame}
	
	\begin{frame}{Q175 --- Model Users}
		\begin{block}{Question}
			Why are model users important in validation?
		\end{block}
		\begin{enumerate}
			\item They build models
			\item They provide feedback on interpretation and usability
			\item They audit
			\item They approve
		\end{enumerate}
		\begin{exampleblock}{Answer: B}\end{exampleblock}
		\begin{block}{Explanation}
			Users understand practical use and can flag issues.
		\end{block}
	\end{frame}
	
	\begin{frame}{Q176 --- Board Responsibilities}
		\begin{block}{Question}
			The board is responsible for:
		\end{block}
		\begin{enumerate}
			\item Building models
			\item Approving the model risk framework and appetite
			\item Auditing models
			\item Using models
		\end{enumerate}
		\begin{exampleblock}{Answer: B}\end{exampleblock}
		\begin{block}{Explanation}
			Board approves framework and sets risk appetite.
		\end{block}
	\end{frame}
	
	\begin{frame}{Q177 --- Model Implementation Tasks}
		\begin{block}{Question}
			Which is NOT a model implementation task?
		\end{block}
		\begin{enumerate}
			\item Model design
			\item Core implementation
			\item System implementation
			\item Independent validation
		\end{enumerate}
		\begin{exampleblock}{Answer: D}\end{exampleblock}
		\begin{block}{Explanation}
			Validation is separate from implementation.
		\end{block}
	\end{frame}
	
	\begin{frame}{Q178 --- Model Adaptation}
		\begin{block}{Question}
			Which scenario requires core adaptation?
		\end{block}
		\begin{enumerate}
			\item Adding a new data interface
			\item Extending code to handle new asset class with mostly reused logic
			\item Building an entirely new system
			\item Adding documentation
		\end{enumerate}
		\begin{exampleblock}{Answer: B}\end{exampleblock}
		\begin{block}{Explanation}
			Core adaptation = modify model core while reusing most code.
		\end{block}
	\end{frame}
	
	\begin{frame}{Q179 --- Data Governance Committee}
		\begin{block}{Question}
			What is the role of a data governance committee?
		\end{block}
		\begin{enumerate}
			\item Operational data entry
			\item Oversight of data policies and practices
			\item Writing code
			\item Building models
		\end{enumerate}
		\begin{exampleblock}{Answer: B}\end{exampleblock}
		\begin{block}{Explanation}
			Committee provides oversight, not operational work.
		\end{block}
	\end{frame}
	
	\begin{frame}{Q180 --- Data Quality}
		\begin{block}{Question}
			Which is NOT a dimension of data quality?
		\end{block}
		\begin{enumerate}
			\item Accuracy
			\item Consistency
			\item Integrity
			\item Profitability
		\end{enumerate}
		\begin{exampleblock}{Answer: D}\end{exampleblock}
		\begin{block}{Explanation}
			Data quality = accuracy, consistency, integrity, completeness.
		\end{block}
	\end{frame}
	
	\begin{frame}{Q181 --- Data Provenance Example}
		\begin{block}{Question}
			A model uses scraped website data without permission. This violates:
		\end{block}
		\begin{enumerate}
			\item Data quality
			\item Data provenance
			\item Data classification
			\item Metadata management
		\end{enumerate}
		\begin{exampleblock}{Answer: B}\end{exampleblock}
		\begin{block}{Explanation}
			Provenance = legal right to use data.
		\end{block}
	\end{frame}
	
	\begin{frame}{Q182 --- Metadata Types}
		\begin{block}{Question}
			Which is NOT a type of metadata?
		\end{block}
		\begin{enumerate}
			\item Descriptive
			\item Structural
			\item Administrative
			\item Predictive
		\end{enumerate}
		\begin{exampleblock}{Answer: D}\end{exampleblock}
		\begin{block}{Explanation}
			Metadata types: descriptive, structural, administrative, reference, statistical.
		\end{block}
	\end{frame}
	
	\begin{frame}{Q183 --- Data Access Controls}
		\begin{block}{Question}
			Which access control model assigns permissions based on job roles?
		\end{block}
		\begin{enumerate}
			\item RBAC
			\item ABAC
			\item MAC
			\item DAC
		\end{enumerate}
		\begin{exampleblock}{Answer: A}\end{exampleblock}
		\begin{block}{Explanation}
			RBAC = Role-Based Access Control.
		\end{block}
	\end{frame}
	
	\begin{frame}{Q184 --- GDPR Principles}
		\begin{block}{Question}
			Which is a core principle of GDPR?
			\begin{enumerate}
				\item Data minimization
				\item Consent
				\item Right to erasure
				\item All of the above
			\end{enumerate}
		\end{block}
		\begin{exampleblock}{Answer: D}\end{exampleblock}
		\begin{block}{Explanation}
			GDPR includes consent, minimization, erasure, and more.
		\end{block}
	\end{frame}
	
	\begin{frame}{Q185 --- Model Governance Framework}
		\begin{block}{Question}
			What is the first step in establishing a model governance framework?
		\end{block}
		\begin{enumerate}
			\item Build a new framework from scratch
			\item Start from existing frameworks and add AI-specific elements
			\item Hire external consultants
			\item Ignore existing frameworks
		\end{enumerate}
		\begin{exampleblock}{Answer: B}\end{exampleblock}
		\begin{block}{Explanation}
			Extend existing frameworks with AI/ML-specific considerations.
		\end{block}
	\end{frame}
	
	\begin{frame}{Q186 --- BIS Survey}
		\begin{block}{Question}
			The BIS survey on AI governance found central banks focus on:
			\begin{enumerate}
				\item Data governance, organization, rules, risk management
				\item Only technology
				\item Only regulation
				\item Only data
			\end{enumerate}
		\end{block}
		\begin{exampleblock}{Answer: A}\end{exampleblock}
		\begin{block}{Explanation}
			Four focus areas: data governance, organization, rules, risk management.
		\end{block}
	\end{frame}
	
	\begin{frame}{Q187 --- GenAI Governance}
		\begin{block}{Question}
			Why is GenAI governance harder than traditional ML governance?
			\begin{enumerate}
				\item Outputs are non-deterministic
				\item Models are opaque
				\item Providers may update models without notice
				\item All of the above
			\end{enumerate}
		\end{block}
		\begin{exampleblock}{Answer: D}\end{exampleblock}
		\begin{block}{Explanation}
			GenAI adds non-determinism, opacity, provider dependence, and rapid change.
		\end{block}
	\end{frame}
	
	\begin{frame}{Q188 --- GenAI Risk}
		\begin{block}{Question}
			Which is NOT a GenAI-specific risk category?
			\begin{enumerate}
				\item Functional risks
				\item Data and information risks
				\item Organizational risks
				\item Weather risks
			\end{enumerate}
			\begin{exampleblock}{Answer: D}\end{exampleblock}
			\begin{block}{Explanation}
				GenAI risk categories: functional, data, organizational, strategic, societal.
			\end{block}
		\end{frame}
		
		\begin{frame}{Q189 --- Model Risk Appetite}
			\begin{block}{Question}
				Model risk appetite should align with:
				\begin{enumerate}
					\item Only regulatory requirements
					\item Overall firm risk appetite and business strategy
					\item Only IT strategy
					\item Only audit findings
				\end{enumerate}
			\end{block}
			\begin{exampleblock}{Answer: B}\end{exampleblock}
			\begin{block}{Explanation}
				Appetite aligns with overall risk appetite and strategy.
			\end{block}
		\end{frame}
		
		\begin{frame}{Q190 --- Internal Audit}
			\begin{block}{Question}
				The role of internal audit in model governance is to:
				\begin{enumerate}
					\item Build models
					\item Independently check the validation process
					\item Set risk appetite
					\item Use models
				\end{enumerate}
			\end{block}
			\begin{exampleblock}{Answer: B}\end{exampleblock}
			\begin{block}{Explanation}
				Internal audit independently checks the validation framework and process.
			\end{block}
		\end{frame}
		
		\begin{frame}{Q191 --- External Audit}
			\begin{block}{Question}
				What role do external auditors play in model risk management?
				\begin{enumerate}
					\item Build models
					\item Independent external check on validation
					\item Set risk appetite
					\item Own models
				\end{enumerate}
			\end{block}
			\begin{exampleblock}{Answer: B}\end{exampleblock}
			\begin{block}{Explanation}
				External auditors are independent external checks.
			\end{block}
		\end{frame}
		
		\begin{frame}{Q192 --- Data Retention}
			\begin{block}{Question}
				A data retention policy should:
				\begin{enumerate}
					\item Keep data forever
					\item Define how long data is retained and when it is deleted
					\item Only apply to structured data
					\item Ignore regulatory requirements
				\end{enumerate}
			\end{block}
			\begin{exampleblock}{Answer: B}\end{exampleblock}
			\begin{block}{Explanation}
				Retention policies define retention periods and deletion schedules.
			\end{block}
		\end{frame}
		
		\begin{frame}{Q193 --- Data Encryption}
			\begin{block}{Question}
				Why is encryption important for data governance?
				\begin{enumerate}
					\item Increases volume
					\item Ensures confidentiality and integrity
					\item Speeds up training
					\item Reduces features
				\end{enumerate}
			\end{block}
			\begin{exampleblock}{Answer: B}\end{exampleblock}
			\begin{block}{Explanation}
				Encryption protects confidentiality and integrity.
			\end{block}
		\end{frame}
		
		\begin{frame}{Q194 --- Gramm-Leach-Bliley}
			\begin{block}{Question}
				The Gramm-Leach-Bliley Act requires financial institutions to:
				\begin{enumerate}
					\item Disclose privacy practices and safeguard customer information
					\item Build AI models
					\item Use only structured data
					\item Ignore state laws
				\end{enumerate}
			\end{block}
			\begin{exampleblock}{Answer: A}\end{exampleblock}
			\begin{block}{Explanation}
				GLBA covers privacy notices and security safeguards.
			\end{block}
		\end{frame}
		
		\begin{frame}{Q195 --- Privacy Act 1974}
			\begin{block}{Question}
				The Privacy Act of 1974 governs:
				\begin{enumerate}
					\item Corporate privacy
					\item Federal agencies' collection and use of personal data
					\item State privacy
					\item International privacy
				\end{enumerate}
				\begin{exampleblock}{Answer: B}\end{exampleblock}
				\begin{block}{Explanation}
					Applies to U.S. federal agencies' records about individuals.
				\end{block}
			\end{frame}
			
			\begin{frame}{Q196 --- Data Stewards}
				\begin{block}{Question}
					Who are data stewards?
					\begin{enumerate}
						\item External consultants
						\item Staff responsible for managing data assets within a governance framework
						\item Software vendors
						\item Regulators
					\end{enumerate}
					\begin{exampleblock}{Answer: B}\end{exampleblock}
					\begin{block}{Explanation}
						Data stewards manage data quality, access, and usage.
					\end{block}
				\end{frame}
				
				\begin{frame}{Q197 --- Model Inventory Contents}
					\begin{block}{Question}
						Which is NOT typically in a model inventory?
						\begin{enumerate}
							\item Model owner
							\item Risk classification
							\item Validation status
							\item Share price
						\end{enumerate}
						\begin{exampleblock}{Answer: D}\end{exampleblock}
						\begin{block}{Explanation}
							Share price is not relevant to model inventory.
						\end{block}
					\end{frame}
					
					\begin{frame}{Q198 --- Model Documentation}
						\begin{block}{Question}
							Which is NOT a typical documentation element?
							\begin{enumerate}
								\item Purpose and intended use
								\item Methodology and assumptions
								\item Marketing plan
								\item Limitations
							\end{enumerate}
							\begin{exampleblock}{Answer: C}\end{exampleblock}
							\begin{block}{Explanation}
								Documentation covers purpose, methodology, limitations, usage.
							\end{block}
						\end{frame}
						
						\begin{frame}{Q199 --- Model Validation Timing}
							\begin{block}{Question}
								Model validation should occur:
								\begin{enumerate}
									\item Only at deployment
									\item Before deployment, periodically, and after material changes
									\item Only after problems arise
									\item Only annually
								\end{enumerate}
								\begin{exampleblock}{Answer: B}\end{exampleblock}
								\begin{block}{Explanation}
									Validation is continuous across the model lifecycle.
								\end{block}
							\end{frame}
							
							\begin{frame}{Q200 --- Model Retirement Process}
								\begin{block}{Question}
									When retiring a model, you should:
									\begin{enumerate}
										\item Just delete it
										\item Document the reason, archive documentation, update inventory
										\item Ignore it
										\item Keep it in production
									\end{enumerate}
									\begin{exampleblock}{Answer: B}\end{exampleblock}
									\begin{block}{Explanation}
										Retirement follows a formal process with documentation.
									\end{block}
								\end{frame}
								
								% =================================================================
								\section{Additional Practice (Q201--Q300)}
								% =================================================================
								
								\begin{frame}{Q201 --- Feature Importance}
									\begin{block}{Question}
										Feature importance scores in XAI provide:
										\begin{enumerate}
											\item Per-prediction attribution
											\item Global ranking of feature influence
											\item Local surrogate models
											\item Privacy preservation
										\end{enumerate}
									\end{block}
									\begin{exampleblock}{Answer: B}\end{exampleblock}
									\begin{block}{Explanation}
										Feature importance = global ranking. SHAP gives per-prediction attribution.
									\end{block}
								\end{frame}
								
								\begin{frame}{Q202 --- LUCID}
									\begin{block}{Question}
										LUCID works by:
										\begin{enumerate}
											\item Predicting outputs from inputs
											\item Working backward from output to input distribution
											\item Using decision trees
											\item Clustering data
										\end{enumerate}
									\end{block}
									\begin{exampleblock}{Answer: B}\end{exampleblock}
									\begin{block}{Explanation}
										LUCID = Locating Unfairness through Canonical Inverse Design.
									\end{block}
								\end{frame}
								
								\begin{frame}{Q203 --- Surrogate Models}
									\begin{block}{Question}
										What is a surrogate model?
										\begin{enumerate}
											\item A simpler interpretable model approximating a complex one
											\item A duplicate model
											\item An encrypted model
											\item A backup model
										\end{enumerate}
									\end{block}
									\begin{exampleblock}{Answer: A}\end{exampleblock}
									\begin{block}{Explanation}
										Surrogate models approximate complex models for interpretability.
									\end{block}
								\end{frame}
								
								\begin{frame}{Q204 --- XAI Challenge}
									\begin{block}{Question}
										What is a key challenge of XAI techniques?
										\begin{enumerate}
											\item They improve accuracy
											\item They are computationally intensive
											\item They eliminate bias
											\item They guarantee fairness
										\end{enumerate}
									\end{block}
									\begin{exampleblock}{Answer: B}\end{exampleblock}
									\begin{block}{Explanation}
										XAI methods can be computationally expensive.
									\end{block}
								\end{frame}
								
								\begin{frame}{Q205 --- Opaqueness Types}
									\begin{block}{Question}
										Which is a form of opaqueness where the individual is unaware of being profiled?
										\begin{enumerate}
											\item Secrecy
											\item Confidentiality
											\item Specialised knowledge
											\item Transparency
										\end{enumerate}
									\end{block}
									\begin{exampleblock}{Answer: A}\end{exampleblock}
									\begin{block}{Explanation}
										Secrecy = unaware of algorithmic decision-making.
									\end{block}
								\end{frame}
								
								\begin{frame}{Q206 --- Confidentiality Opaqueness}
									\begin{block}{Question}
										Confidentiality opaqueness refers to:
										\begin{enumerate}
											\item Being unaware of algorithm existence
											\item Lack of access to algorithm internals
											\item Lack of specialised knowledge
											\item Full transparency
										\end{enumerate}
									\end{block}
									\begin{exampleblock}{Answer: B}\end{exampleblock}
									\begin{block}{Explanation}
										Confidentiality = aware of decision-making but not internals.
									\end{block}
								\end{frame}
								
								\begin{frame}{Q207 --- Specialised Knowledge}
									\begin{block}{Question}
										Specialised knowledge opaqueness means:
										\begin{enumerate}
											\item No one can understand
											\item Understanding requires expertise
											\item Legal restriction
											\item System error
										\end{enumerate}
									\end{block}
									\begin{exampleblock}{Answer: B}\end{exampleblock}
									\begin{block}{Explanation}
										Even with access, expertise is required to interpret.
									\end{block}
								\end{frame}
								
								\begin{frame}{Q208 --- Recursive Validation}
									\begin{block}{Question}
										Recursive validation means:
										\begin{enumerate}
											\item Validation is repeated
											\item Validation itself should be critically examined
											\item Validation is automated
											\item Validation is optional
										\end{enumerate}
									\end{block}
									\begin{exampleblock}{Answer: B}\end{exampleblock}
									\begin{block}{Explanation}
										Validation should not be exempt from critical examination.
									\end{block}
								\end{frame}
								
								\begin{frame}{Q209 --- No Valid Model}
									\begin{block}{Question}
										Why is there no perfectly valid model?
										\begin{enumerate}
											\item Models are always wrong
											\item Models simplify reality under assumptions
											\item Data is always biased
											\item Regulators forbid it
										\end{enumerate}
									\end{block}
									\begin{exampleblock}{Answer: B}\end{exampleblock}
									\begin{block}{Explanation}
										All models are simplifications; the goal is to withstand invalidation attempts.
									\end{block}
								\end{frame}
								
								\begin{frame}{Q210 --- Effective Challenge}
									\begin{block}{Question}
										Effective challenge requires:
										\begin{enumerate}
											\item Independence, competence, and influence
											\item Only technical skills
											\item Only independence
											\item Only seniority
										\end{enumerate}
									\end{block}
									\begin{exampleblock}{Answer: A}\end{exampleblock}
									\begin{block}{Explanation}
										Effective challenge needs independence, competence, and influence.
									\end{block}
								\end{frame}
								
								\begin{frame}{Q211 --- GenAI Validation}
									\begin{block}{Question}
										Which is a GenAI-specific validation challenge?
										\begin{enumerate}
											\item Open-ended output
											\item Hallucinations
											\item Multi-modal inputs
											\item All of the above
										\end{enumerate}
									\end{block}
									\begin{exampleblock}{Answer: D}\end{exampleblock}
									\begin{block}{Explanation}
										GenAI validation challenges include open-ended output, hallucinations, multi-modality.
									\end{block}
								\end{frame}
								
								\begin{frame}{Q212 --- Model Design Motivation}
									\begin{block}{Question}
										Which is NOT a common motivation for new model development?
										\begin{enumerate}
											\item Changing market conditions
											\item Regulatory pressure
											\item Recent crisis losses
											\item Decrease in staff
										\end{enumerate}
									\end{block}
									\begin{exampleblock}{Answer: D}\end{exampleblock}
									\begin{block}{Explanation}
										Staffing changes are not typically drivers for model development.
									\end{block}
								\end{frame}
								
								\begin{frame}{Q213 --- Quants Background}
									\begin{block}{Question}
										Model developers typically have backgrounds in:
										\begin{enumerate}
											\item Mathematics, physics, engineering, or computer science
											\item Only finance
											\item Only economics
											\item Only marketing
										\end{enumerate}
									\end{block}
									\begin{exampleblock}{Answer: A}\end{exampleblock}
									\begin{block}{Explanation}
										Quants have quantitative scientific backgrounds.
									\end{block}
								\end{frame}
								
								\begin{frame}{Q214 --- Discretization}
									\begin{block}{Question}
										Discretization in modeling is used to:
										\begin{enumerate}
											\item Increase complexity
											\item Approximate continuous problems on a finite grid
											\item Eliminate assumptions
											\item Ignore boundary conditions
										\end{enumerate}
									\end{block}
									\begin{exampleblock}{Answer: B}\end{exampleblock}
									\begin{block}{Explanation}
										Discretization approximates continuous problems numerically.
									\end{block}
								\end{frame}
								
								\begin{frame}{Q215 --- Approximations}
									\begin{block}{Question}
										Square-root-of-time scaling for VaR assumes:
										\begin{enumerate}
											\item Arbitrary scaling
											\item Normal distribution and temporal independence
											\item Deterministic volatility
											\item Perfect correlation
										\end{enumerate}
									\end{block}
									\begin{exampleblock}{Answer: B}\end{exampleblock}
									\begin{block}{Explanation}
										Requires normality and independence, often violated in practice.
									\end{block}
								\end{frame}
								
								\begin{frame}{Q216 --- Numerical Evaluation}
									\begin{block}{Question}
										Numerical evaluation differs from approximation by:
										\begin{enumerate}
											\item Lower precision
											\item Arbitrary precision potential
											\item No parameters
											\item No convergence
										\end{enumerate}
									\end{block}
									\begin{exampleblock}{Answer: B}\end{exampleblock}
									\begin{block}{Explanation}
										Numerical methods can reach arbitrary precision.
									\end{block}
								\end{frame}
								
								\begin{frame}{Q217 --- Random Numbers}
									\begin{block}{Question}
										Computers generate random numbers using:
										\begin{enumerate}
											\item True randomness
											\item Pseudo-random number generators (RNGs)
											\item Physical dice
											\item Manual selection
										\end{enumerate}
									\end{block}
									\begin{exampleblock}{Answer: B}\end{exampleblock}
									\begin{block}{Explanation}
										Computers use deterministic pseudo-RNGs seeded initially.
									\end{block}
								\end{frame}
								
								\begin{frame}{Q218 --- Model Implementation}
									\begin{block}{Question}
										Which is a model implementation task?
										\begin{enumerate}
											\item Building user interfaces
											\item Model design
											\item Independent validation
											\item Data cleaning
										\end{enumerate}
									\end{block}
									\begin{exampleblock}{Answer: B}\end{exampleblock}
									\begin{block}{Explanation}
										Model design is part of implementation (design → core → system).
									\end{block}
								\end{frame}
								
								\begin{frame}{Q219 --- System Implementation}
									\begin{block}{Question}
										System implementation involves:
										\begin{enumerate}
											\item Only the model core
											\item Integration with UI, data inputs, scheduling, outputs
											\item Only documentation
											\item Only validation
										\end{enumerate}
									\end{block}
									\begin{exampleblock}{Answer: B}\end{exampleblock}
									\begin{block}{Explanation}
										System implementation wraps the core with UI, inputs, scheduling.
									\end{block}
								\end{frame}
								
								\begin{frame}{Q220 --- Model Adaptation}
									\begin{block}{Question}
										Core adaptation differs from system adaptation by:
										\begin{enumerate}
											\item Adapting the model core vs the surrounding system
											\item Being easier
											\item Being cheaper
											\item Being optional
										\end{enumerate}
									\end{block}
									\begin{exampleblock}{Answer: A}\end{exampleblock}
									\begin{block}{Explanation}
										Core adaptation = model core changes. System adaptation = surrounding system.
									\end{block}
								\end{frame}
								
								\begin{frame}{Q221 --- Adapting Large Systems}
									\begin{block}{Question}
										Adapting large or complex systems is problematic because:
										\begin{enumerate}
											\item They are always wrong
											\item Original developers may have left; documentation may be incomplete
											\item They cannot be adapted
											\item Regulators forbid changes
										\end{enumerate}
									\end{block}
									\begin{exampleblock}{Answer: B}\end{exampleblock}
									\begin{block}{Explanation}
										Lost knowledge, incomplete documentation make adaptation risky.
									\end{block}
								\end{frame}
								
								\begin{frame}{Q222 --- Running a Model}
									\begin{block}{Question}
										What is a key challenge when running a model?
										\begin{enumerate}
											\item Data gathering and processing
											\item Scheduling and interpreting results
											\item Both A and B
											\item None of the above
										\end{enumerate}
									\end{block}
									\begin{exampleblock}{Answer: C}\end{exampleblock}
									\begin{block}{Explanation}
										Running a model involves data gathering, scheduling, and interpretation.
									\end{block}
								\end{frame}
								
								\begin{frame}{Q223 --- Scheduled vs Ad Hoc Runs}
									\begin{block}{Question}
										Scheduled runs differ from ad hoc runs by:
										\begin{enumerate}
											\item No manual intervention; regular schedule
											\item Requiring manual inputs
											\item Being optional
											\item Never producing output
										\end{enumerate}
									\end{block}
									\begin{exampleblock}{Answer: A}\end{exampleblock}
									\begin{block}{Explanation}
										Scheduled = automated; ad hoc = manual, on-demand.
									\end{block}
								\end{frame}
								
								\begin{frame}{Q224 --- Misinterpretation Causes}
									\begin{block}{Question}
										Which is a cause of model misinterpretation?
										\begin{enumerate}
											\item Overemphasis on point estimates
											\item Ignoring qualitative factors
											\item Confirmation bias
											\item All of the above
										\end{enumerate}
									\end{block}
									\begin{exampleblock}{Answer: D}\end{exampleblock}
									\begin{block}{Explanation}
										Common causes: point-estimate focus, ignoring context, confirmation bias.
									\end{block}
								\end{frame}
								
								\begin{frame}{Q225 --- Model Risk Ranking}
									\begin{block}{Question}
										Model risk ranking typically considers:
										\begin{enumerate}
											\item Materiality, complexity, exposure
											\item Only materiality
											\item Only complexity
											\item Only cost
										\end{enumerate}
									\end{block}
									\begin{exampleblock}{Answer: A}\end{exampleblock}
									\begin{block}{Explanation}
										Risk ranking = materiality × complexity × exposure.
									\end{block}
								\end{frame}
								
								\begin{frame}{Q226 --- Model Performance Metrics}
									\begin{block}{Question}
										For classification models, which metrics are commonly used?
										\begin{enumerate}
											\item Accuracy, Precision, Recall, F1
											\item MAE, MSE
											\item Only accuracy
											\item Only precision
										\end{enumerate}
									\end{block}
									\begin{exampleblock}{Answer: A}\end{exampleblock}
									\begin{block}{Explanation}
										Classification: accuracy, precision, recall, F1, AUC.
									\end{block}
								\end{frame}
								
								\begin{frame}{Q227 --- Regression Metrics}
									\begin{block}{Question}
										For regression models, which metrics are commonly used?
										\begin{enumerate}
											\item Accuracy, Precision
											\item MAE, MSE, RMSE, R$^2$
											\item Recall, F1
											\item AUC
										\end{enumerate}
									\end{block}
									\begin{exampleblock}{Answer: B}\end{exampleblock}
									\begin{block}{Explanation}
										Regression: MAE, MSE, RMSE, R$^2$.
									\end{block}
								\end{frame}
								
								\begin{frame}{Q228 --- GenAI Recalibration}
									\begin{block}{Question}
										GenAI models require:
										\begin{enumerate}
											\item No recalibration
											\item More frequent recalibration than traditional models
											\item Only annual recalibration
											\item Only initial calibration
										\end{enumerate}
									\end{block}
									\begin{exampleblock}{Answer: B}\end{exampleblock}
									\begin{block}{Explanation}
										Rapid change requires more frequent recalibration.
									\end{block}
								\end{frame}
								
								\begin{frame}{Q229 --- Model Use Limits}
									\begin{block}{Question}
										Model use limits should be documented in:
										\begin{enumerate}
											\item Marketing materials
											\item The model inventory
											\item The code
											\item The audit report
										\end{enumerate}
									\end{block}
									\begin{exampleblock}{Answer: B}\end{exampleblock}
									\begin{block}{Explanation}
										Use limits are documented in the model inventory.
									\end{block}
								\end{frame}
								
								\begin{frame}{Q230 --- Model Risk Culture Guidelines}
									\begin{block}{Question}
										Which is NOT a guideline for a strong modeling culture?
										\begin{enumerate}
											\item Awareness
											\item Transparency
											\item Experience
											\item Secrecy
										\end{enumerate}
									\end{block}
									\begin{exampleblock}{Answer: D}\end{exampleblock}
									\begin{block}{Explanation}
										Culture guidelines: awareness, transparency, experience.
									\end{block}
								\end{frame}
								
								\begin{frame}{Q231 --- Change Management Process}
									\begin{block}{Question}
										What is included in change management?
										\begin{enumerate}
											\item Change request
											\item Impact assessment
											\item Documentation and validation
											\item All of the above
										\end{enumerate}
									\end{block}
									\begin{exampleblock}{Answer: D}\end{exampleblock}
									\begin{block}{Explanation}
										Change management includes requests, impact, documentation, validation.
									\end{block}
								\end{frame}
								
								\begin{frame}{Q232 --- Independent Review}
									\begin{block}{Question}
										Independent review of model changes is performed by:
										\begin{enumerate}
											\item The original developers
											\item Parties not involved in development
											\item External vendors only
											\item Regulators
										\end{enumerate}
									\end{block}
									\begin{exampleblock}{Answer: B}\end{exampleblock}
									\begin{block}{Explanation}
										Review must be independent from development.
									\end{block}
								\end{frame}
								
								\begin{frame}{Q233 --- Approval}
									\begin{block}{Question}
										Model change approval is obtained through:
										\begin{enumerate}
											\item Informal discussion
											\item Defined governance process
											\item Only IT approval
											\item Only the CRO
										\end{enumerate}
									\end{block}
									\begin{exampleblock}{Answer: B}\end{exampleblock}
									\begin{block}{Explanation}
										Approvals go through defined governance.
									\end{block}
								\end{frame}
								
								\begin{frame}{Q234 --- Implementation and Monitoring}
									\begin{block}{Question}
										After implementation, model changes are:
										\begin{enumerate}
											\item Ignored
											\item Monitored to ensure objectives are met
											\item Never changed
											\item Retired immediately
										\end{enumerate}
									\end{block}
									\begin{exampleblock}{Answer: B}\end{exampleblock}
									\begin{block}{Explanation}
										Post-implementation monitoring ensures objectives are met.
									\end{block}
								\end{frame}
								
								\begin{frame}{Q235 --- AI/ML vs Traditional Review Frequency}
									\begin{block}{Question}
										AI/ML models require more frequent monitoring than traditional models because:
										\begin{enumerate}
											\item They are cheaper
											\item Higher complexity and susceptibility to data drift
											\item Regulators demand it
											\item All of the above
										\end{enumerate}
									\end{block}
									\begin{exampleblock}{Answer: B}\end{exampleblock}
									\begin{block}{Explanation}
										Complexity and data drift drive more frequent review.
									\end{block}
								\end{frame}
								
								\begin{frame}{Q236 --- Validation Stakeholders}
									\begin{block}{Question}
										Which stakeholder has expanded role in AI/ML validation?
										\begin{enumerate}
											\item Only IT
											\item HR, compliance, operational risk, and others
											\item Only model developers
											\item Only regulators
										\end{enumerate}
									\end{block}
									\begin{exampleblock}{Answer: B}\end{exampleblock}
									\begin{block}{Explanation}
										AI/ML validation involves expanded stakeholder set.
									\end{block}
								\end{frame}
								
								\begin{frame}{Q237 --- Explainability Focus}
									\begin{block}{Question}
										For black-box models, validation should focus on:
										\begin{enumerate}
											\item Code review only
											\item Explainability, benchmarking, and output testing
											\item Ignoring outputs
											\item Only documentation
										\end{enumerate}
									\end{block}
									\begin{exampleblock}{Answer: B}\end{exampleblock}
									\begin{block}{Explanation}
										Black-box models require output-focused testing and explainability.
									\end{block}
								\end{frame}
								
								\begin{frame}{Q238 --- Benchmarking}
									\begin{block}{Question}
										Benchmarking AI/ML models against classical models is recommended to:
										\begin{enumerate}
											\item Replace classical models
											\item Understand reasons for deviations
											\item Ignore classical models
											\item Simplify validation
										\end{enumerate}
									\end{block}
									\begin{exampleblock}{Answer: B}\end{exampleblock}
									\begin{block}{Explanation}
										Benchmarking helps understand deviations.
									\end{block}
								\end{frame}
								
								\begin{frame}{Q239 --- K-Fold Cross-Validation for AI}
									\begin{block}{Question}
										K-fold cross-validation is suitable when:
										\begin{enumerate}
											\item Data is time-series
											\item Data has no natural ordering
											\item Data is small
											\item Both B and C
										\end{enumerate}
									\end{block}
									\begin{exampleblock}{Answer: D}\end{exampleblock}
									\begin{block}{Explanation}
										K-fold CV works well for small unordered data.
									\end{block}
								\end{frame}
								
								\begin{frame}{Q240 --- Backtesting}
									\begin{block}{Question}
										Traditional backtesting methods may become ineffective for AI/ML because:
										\begin{enumerate}
											\item No inputs
											\item No straightforward relation between inputs and outputs
											\item Data is too small
											\item Models are too simple
										\end{enumerate}
									\end{block}
									\begin{exampleblock}{Answer: B}\end{exampleblock}
									\begin{block}{Explanation}
										Nonlinear ML models complicate traditional backtesting.
									\end{block}
								\end{frame}
								
								\begin{frame}{Q241 --- Reporting Component}
									\begin{block}{Question}
										Model output reporting should verify:
										\begin{enumerate}
											\item Only numbers
											\item Accuracy, completeness, appropriateness, and limitations
											\item Only speed
											\item Only cost
										\end{enumerate}
									\end{block}
									\begin{exampleblock}{Answer: B}\end{exampleblock}
									\begin{block}{Explanation}
										Reports verify accuracy, completeness, appropriateness, limitations.
									\end{block}
								\end{frame}
								
								\begin{frame}{Q242 --- Model Validation Definition}
									\begin{block}{Question}
										Model validation is defined as:
										\begin{enumerate}
											\item Building the model
											\item Rigorously testing model, inputs, outputs, governance
											\item Using the model
											\item Marketing the model
										\end{enumerate}
									\end{block}
									\begin{exampleblock}{Answer: B}\end{exampleblock}
									\begin{block}{Explanation}
										Validation tests model, inputs, outputs, governance.
									\end{block}
								\end{frame}
								
								\begin{frame}{Q243 --- Validation Lifecycle}
									\begin{block}{Question}
										Model validation lifecycle starts with:
										\begin{enumerate}
											\item Retirement
											\item Identification and inventory
											\item Implementation
											\item Usage
										\end{enumerate}
									\end{block}
									\begin{exampleblock}{Answer: B}\end{exampleblock}
									\begin{block}{Explanation}
										Lifecycle: identification → inventory → validation → production → retirement.
									\end{block}
								\end{frame}
								
								\begin{frame}{Q244 --- Initial Validation Goals}
									\begin{block}{Question}
										Initial validation ensures:
										\begin{enumerate}
											\item Operational feasibility
											\item Proper documentation
											\item User training
											\item All of the above
										\end{enumerate}
									\end{block}
									\begin{exampleblock}{Answer: D}\end{exampleblock}
									\begin{block}{Explanation}
										Initial validation covers feasibility, documentation, and training.
									\end{block}
								\end{frame}
								
								\begin{frame}{Q245 --- Ongoing Validation}
									\begin{block}{Question}
										Ongoing validation confirms:
										\begin{enumerate}
											\item Model remains aligned with intended purpose
											\item Assumptions remain valid
											\item Data still appropriate
											\item All of the above
										\end{enumerate}
									\end{block}
									\begin{exampleblock}{Answer: D}\end{exampleblock}
									\begin{block}{Explanation}
										Ongoing validation confirms alignment, assumptions, data, performance.
									\end{block}
								\end{frame}
								
								\begin{frame}{Q246 --- Validation Frequency}
									\begin{block}{Question}
										Validation frequency depends on:
										\begin{enumerate}
											\item Model risk tier
											\item Calendar year
											\item Only regulatory schedule
											\item Model age
										\end{enumerate}
									\end{block}
									\begin{exampleblock}{Answer: A}\end{exampleblock}
									\begin{block}{Explanation}
										Higher risk = more frequent validation.
									\end{block}
								\end{frame}
								
								\begin{frame}{Q247 --- Validation Terminology}
									\begin{block}{Question}
										Alternative terms for model validation include:
										\begin{enumerate}
											\item Model review, QA, evaluation, vetting
											\item Model development
											\item Model implementation
											\item Model usage
										\end{enumerate}
									\end{block}
									\begin{exampleblock}{Answer: A}\end{exampleblock}
									\begin{block}{Explanation}
										Validation = review, QA, evaluation, vetting.
									\end{block}
								\end{frame}
								
								\begin{frame}{Q248 --- Suitability Focus}
									\begin{block}{Question}
										The ultimate goal of validation is to assess:
										\begin{enumerate}
											\item Whether the model is appropriate and effective for its purpose
											\item Whether the model is complex
											\item Whether the model is expensive
											\item Whether the model is simple
										\end{enumerate}
									\end{block}
									\begin{exampleblock}{Answer: A}\end{exampleblock}
									\begin{block}{Explanation}
										Validation focuses on fitness for purpose.
									\end{block}
								\end{frame}
								
								\begin{frame}{Q249 --- Recursive Validation}
									\begin{block}{Question}
										Recursive validation requires:
										\begin{enumerate}
											\item Validation not exempt from critical examination
											\item Validation is automatic
											\item Validation is optional
											\item Validation is one-time
										\end{enumerate}
									\end{block}
									\begin{exampleblock}{Answer: A}\end{exampleblock}
									\begin{block}{Explanation}
										Validation should itself be critically examined.
									\end{block}
								\end{frame}
								
								\begin{frame}{Q250 --- No Valid Model Concept}
									\begin{block}{Question}
										The idea ``no valid model'' means:
										\begin{enumerate}
											\item Models are useless
											\item Models are simplifications; validation tries to invalidate them
											\item Models are always correct
											\item Models cannot be validated
										\end{enumerate}
									\end{block}
									\begin{exampleblock}{Answer: B}\end{exampleblock}
									\begin{block}{Explanation}
										Models simplify reality; validation tests whether they withstand challenge.
									\end{block}
								\end{frame}
								
								\begin{frame}{Q251 --- Usefulness vs Validity}
									\begin{block}{Question}
										Models can be useful even if:
										\begin{enumerate}
											\item They are perfect
											\item They have identified weaknesses and limitations
											\item They are always wrong
											\item They are not used
										\end{enumerate}
									\end{block}
									\begin{exampleblock}{Answer: B}\end{exampleblock}
									\begin{block}{Explanation}
										Useful models can have limitations if documented and understood.
									\end{block}
								\end{frame}
								
								\begin{frame}{Q252 --- Effective Challenge Definition}
									\begin{block}{Question}
										Effective challenge requires:
										\begin{enumerate}
											\item Independent, competent, influential reviewers
											\item Only seniority
											\item Only technical skill
											\item Only independence
										\end{enumerate}
									\end{block}
									\begin{exampleblock}{Answer: A}\end{exampleblock}
									\begin{block}{Explanation}
										Independence + competence + influence.
									\end{block}
								\end{frame}
								
								\begin{frame}{Q253 --- GenAI Validation Challenges}
									\begin{block}{Question}
										Which is NOT a GenAI validation challenge?
										\begin{enumerate}
											\item Open-ended output
											\item Hallucination
											\item Data drift
											\item Perfect accuracy
										\end{enumerate}
									\end{block}
									\begin{exampleblock}{Answer: D}\end{exampleblock}
									\begin{block}{Explanation}
										Perfect accuracy is not a challenge.
									\end{block}
								\end{frame}
								
								\begin{frame}{Q254 --- Model Design Error}
									\begin{block}{Question}
										Model design errors may only become apparent:
										\begin{enumerate}
											\item During deployment
											\item During future crises
											\item Immediately
											\item Never
										\end{enumerate}
									\end{block}
									\begin{exampleblock}{Answer: B}\end{exampleblock}
									\begin{block}{Explanation}
										Design errors often surface during crises.
									\end{block}
								\end{frame}
								
								\begin{frame}{Q255 --- Gaussian Copula}
									\begin{block}{Question}
										The Gaussian copula model was criticised for:
										\begin{enumerate}
											\item Underestimating risk in CDOs
											\item Overestimating risk
											\item Being too complex
											\item Being too simple
										\end{enumerate}
									\end{block}
									\begin{exampleblock}{Answer: A}\end{exampleblock}
									\begin{block}{Explanation}
										Gaussian copula contributed to CDO mispricing pre-2008.
									\end{block}
								\end{frame}
								
								\begin{frame}{Q256 --- Black-Scholes}
									\begin{block}{Question}
										Black-Scholes model assumes:
										\begin{enumerate}
											\item Constant volatility
											\item Stochastic volatility
											\item Jump diffusion
											\item None of the above
										\end{enumerate}
									\end{block}
									\begin{exampleblock}{Answer: A}\end{exampleblock}
									\begin{block}{Explanation}
										Constant volatility assumption criticised in practice.
									\end{block}
								\end{frame}
								
								\begin{frame}{Q257 --- Balance in Model Design}
									\begin{block}{Question}
										Striking a balance in model design means:
										\begin{enumerate}
											\item Maximum complexity
											\item Robustness vs accounting for shortcomings
											\item Ignoring limitations
											\item Maximum simplicity
										\end{enumerate}
									\end{block}
									\begin{exampleblock}{Answer: B}\end{exampleblock}
									\begin{block}{Explanation}
										Balance robustness and known limitations.
									\end{block}
								\end{frame}
								
								\begin{frame}{Q258 --- Discretization Error}
									\begin{block}{Question}
										Discretization error arises from:
										\begin{enumerate}
											\item Approximating continuous problems on a grid
											\item Ignoring data
											\item Too much data
											\item Perfect precision
										\end{enumerate}
									\end{block}
									\begin{exampleblock}{Answer: A}\end{exampleblock}
									\begin{block}{Explanation}
										Grid approximation introduces discretization error.
									\end{block}
								\end{frame}
								
								\begin{frame}{Q259 --- Stability in Discretization}
									\begin{block}{Question}
										In discretization, stability refers to:
										\begin{enumerate}
											\item Small errors not becoming unbounded
											\item Speed of computation
											\item Memory usage
											\item Model accuracy
										\end{enumerate}
									\end{block}
									\begin{exampleblock}{Answer: A}\end{exampleblock}
									\begin{block}{Explanation}
										Stability = small errors stay small.
									\end{block}
								\end{frame}
								
								\begin{frame}{Q260 --- Approximation Validity}
									\begin{block}{Question}
										Approximations in modeling are valid when:
										\begin{enumerate}
											\item Always
											\item Under certain conditions and understood limitations
											\item Never
											\item Only for simple models
										\end{enumerate}
									\end{block}
									\begin{exampleblock}{Answer: B}\end{exampleblock}
									\begin{block}{Explanation}
										Approximations valid under specific conditions.
									\end{block}
								\end{frame}
								
								\begin{frame}{Q261 --- Numerical Evaluation}
									\begin{block}{Question}
										Numerical evaluation is available for:
										\begin{enumerate}
											\item Only linear problems
											\item Many model components with mathematical justification
											\item Only simple problems
											\item No problems
										\end{enumerate}
									\end{block}
									\begin{exampleblock}{Answer: B}\end{exampleblock}
									\begin{block}{Explanation}
										Numerical methods available broadly with conditions.
									\end{block}
								\end{frame}
								
								\begin{frame}{Q262 --- Floating Point}
									\begin{block}{Question}
										Floating-point arithmetic introduces:
										\begin{enumerate}
											\item Exact results
											\item Rounding errors
											\item No error
											\item Faster training
										\end{enumerate}
									\end{block}
									\begin{exampleblock}{Answer: B}\end{exampleblock}
									\begin{block}{Explanation}
										Floating-point arithmetic causes rounding errors.
									\end{block}
								\end{frame}
								
								\begin{frame}{Q263 --- RNG Quality}
									\begin{block}{Question}
										RNG quality is assessed using:
										\begin{enumerate}
											\item Test suites for randomness
											\item Only speed
											\item Only reproducibility
											\item Only simplicity
										\end{enumerate}
									\end{block}
									\begin{exampleblock}{Answer: A}\end{exampleblock}
									\begin{block}{Explanation}
										Test suites evaluate RNG quality.
									\end{block}
								\end{frame}
								
								\begin{frame}{Q264 --- Model Design Documentation}
									\begin{block}{Question}
										Model design should be documented to:
										\begin{enumerate}
											\item Facilitate implementation
											\item Enable validation
											\item Support audit
											\item All of the above
										\end{enumerate}
									\end{block}
									\begin{exampleblock}{Answer: D}\end{exampleblock}
									\begin{block}{Explanation}
										Documentation supports implementation, validation, audit.
									\end{block}
								\end{frame}
								
								\begin{frame}{Q265 --- Core Implementation}
									\begin{block}{Question}
										Core implementation involves:
										\begin{enumerate}
											\item Building model logic as a black box with interfaces
											\item Building UI
											\item Data collection
											\item Validation
										\end{enumerate}
									\end{block}
									\begin{exampleblock}{Answer: A}\end{exampleblock}
									\begin{block}{Explanation}
										Core = model logic with specified interfaces.
									\end{block}
								\end{frame}
								
								\begin{frame}{Q266 --- System Implementation}
									\begin{block}{Question}
										System implementation provides:
										\begin{enumerate}
											\item Only model core
											\item UI, data inputs, scheduling, outputs
											\item Only documentation
											\item Only validation
										\end{enumerate}
									\end{block}
									\begin{exampleblock}{Answer: B}\end{exampleblock}
									\begin{block}{Explanation}
										System surrounds core with UI, inputs, scheduling, outputs.
									\end{block}
								\end{frame}
								
								\begin{frame}{Q267 --- Model Adaptation Types}
									\begin{block}{Question}
										Which is NOT a model adaptation task?
										\begin{enumerate}
											\item Model adaptation
											\item Core adaptation
											\item System adaptation
											\item Independent validation
										\end{enumerate}
									\end{block}
									\begin{exampleblock}{Answer: D}\end{exampleblock}
									\begin{block}{Explanation}
										Validation is separate.
									\end{block}
								\end{frame}
								
								\begin{frame}{Q268 --- Adaptation Risk}
									\begin{block}{Question}
										Risk in adapting large systems:
										\begin{enumerate}
											\item Introduction of errors into previously functional code
											\item Faster execution
											\item Lower cost
											\item Simplicity
										\end{enumerate}
									\end{block}
									\begin{exampleblock}{Answer: A}\end{exampleblock}
									\begin{block}{Explanation}
										Adaptation can introduce errors in working code.
									\end{block}
								\end{frame}
								
								\begin{frame}{Q269 --- Robustness in Modeling}
									\begin{block}{Question}
										Robustness in modeling means:
										\begin{enumerate}
											\item Models run fast
											\item Models perform well under varied conditions
											\item Models are simple
											\item Models are complex
										\end{enumerate}
									\end{block}
									\begin{exampleblock}{Answer: B}\end{exampleblock}
									\begin{block}{Explanation}
										Robustness = performance across conditions.
									\end{block}
								\end{frame}
								
								\begin{frame}{Q270 --- Scheduled Runs}
									\begin{block}{Question}
										Scheduled runs focus on:
										\begin{enumerate}
											\item IT reliability and safety
											\item Only accuracy
											\item Only cost
											\item Only speed
										\end{enumerate}
									\end{block}
									\begin{exampleblock}{Answer: A}\end{exampleblock}
									\begin{block}{Explanation}
										Scheduled runs prioritize reliability and safety.
									\end{block}
								\end{frame}
								
								\begin{frame}{Q271 --- Ad Hoc Runs}
									\begin{block}{Question}
										Ad hoc runs require:
										\begin{enumerate}
											\item No intervention
											\item Manual intervention and knowledge of interfaces
											\item Only automation
											\item Only documentation
										\end{enumerate}
									\end{block}
									\begin{exampleblock}{Answer: B}\end{exampleblock}
									\begin{block}{Explanation}
										Ad hoc requires manual intervention and interface knowledge.
									\end{block}
								\end{frame}
								
								\begin{frame}{Q272 --- Model Misinterpretation}
									\begin{block}{Question}
										Which is NOT a cause of misinterpretation?
										\begin{enumerate}
											\item Overemphasis on point estimates
											\item Ignoring qualitative factors
											\item Confirmation bias
											\item Too much training data
										\end{enumerate}
									\end{block}
									\begin{exampleblock}{Answer: D}\end{exampleblock}
									\begin{block}{Explanation}
										Too much data is not a misinterpretation cause.
									\end{block}
								\end{frame}
								
								\begin{frame}{Q273 --- Model Risk Ranking}
									\begin{block}{Question}
										Risk ranking helps:
										\begin{enumerate}
											\item Prioritize validation resources
											\item Reduce data
											\item Increase features
											\item Ignore governance
										\end{enumerate}
									\end{block}
									\begin{exampleblock}{Answer: A}\end{exampleblock}
									\begin{block}{Explanation}
										Risk ranking prioritizes resources.
									\end{block}
								\end{frame}
								
								\begin{frame}{Q274 --- Qualitative Factors}
									\begin{block}{Question}
										When interpreting model outputs, qualitative factors:
										\begin{enumerate}
											\item Are irrelevant
											\item Should be considered alongside quantitative results
											\item Replace quantitative
											\item Never matter
										\end{enumerate}
									\end{block}
									\begin{exampleblock}{Answer: B}\end{exampleblock}
									\begin{block}{Explanation}
										Qualitative factors complement quantitative results.
									\end{block}
								\end{frame}
								
								\begin{frame}{Q275 --- Model Use Test}
									\begin{block}{Question}
										The use test is:
										\begin{enumerate}
											\item Quantitative
											\item Qualitative, people-focused
											\item Only numbers
											\item Only automated
										\end{enumerate}
									\end{block}
									\begin{exampleblock}{Answer: B}\end{exampleblock}
									\begin{block}{Explanation}
										Use test = qualitative, people-focused.
									\end{block}
								\end{frame}
								
								\begin{frame}{Q276 --- Use Test Audience}
									\begin{block}{Question}
										Use test results are typically presented to:
										\begin{enumerate}
											\item Junior staff
											\item Senior management
											\item Only auditors
											\item Only regulators
										\end{enumerate}
									\end{block}
									\begin{exampleblock}{Answer: B}\end{exampleblock}
									\begin{block}{Explanation}
										Use test results are for senior management.
									\end{block}
								\end{frame}
								
								\begin{frame}{Q277 --- Validation Report}
									\begin{block}{Question}
										A model validation culminates in:
										\begin{enumerate}
											\item A code file
											\item A validation report with pass/fail rating
											\item Only a database
											\item Nothing
										\end{enumerate}
									\end{block}
									\begin{exampleblock}{Answer: B}\end{exampleblock}
									\begin{block}{Explanation}
										Validation ends with a report and rating.
									\end{block}
								\end{frame}
								
								\begin{frame}{Q278 --- Validation Frequency}
									\begin{block}{Question}
										OCC recommends model performance review at least:
										\begin{enumerate}
											\item Every 5 years
											\item Annually or more frequently if warranted
											\item Never
											\item Biennially
										\end{enumerate}
									\end{block}
									\begin{exampleblock}{Answer: B}\end{exampleblock}
									\begin{block}{Explanation}
										Annual minimum or more frequently.
									\end{block}
								\end{frame}
								
								\begin{frame}{Q279 --- Validation Activities}
									\begin{block}{Question}
										Validation includes:
										\begin{enumerate}
											\item Data source assessment
											\item Model assumptions review
											\item Implementation platform review
											\item All of the above
										\end{enumerate}
									\end{block}
									\begin{exampleblock}{Answer: D}\end{exampleblock}
									\begin{block}{Explanation}
										Validation covers data, assumptions, implementation, more.
									\end{block}
								\end{frame}
								
								\begin{frame}{Q280 --- Backtesting Frequency}
									\begin{block}{Question}
										Backtesting frequency depends on:
										\begin{enumerate}
											\item Model use and risk
											\item Only regulatory schedule
											\item Only model age
											\item Only data size
										\end{enumerate}
									\end{block}
									\begin{exampleblock}{Answer: A}\end{exampleblock}
									\begin{block}{Explanation}
										Backtesting frequency depends on model use and risk.
									\end{block}
								\end{frame}
								
								\begin{frame}{Q281 --- VaR Backtesting}
									\begin{block}{Question}
										VaR models require backtesting at:
										\begin{enumerate}
											\item Annual frequency
											\item Daily frequency
											\item Monthly frequency
											\item Quarterly frequency
										\end{enumerate}
										\begin{exampleblock}{Answer: B}\end{exampleblock}
										\begin{block}{Explanation}
											VaR backtesting is daily.
										\end{block}
									\end{frame}
									
									\begin{frame}{Q282 --- Credit Risk Backtesting}
										\begin{block}{Question}
											Credit risk models may be backtested:
											\begin{enumerate}
												\item Daily
												\item Monthly
												\item Annually
												\item Never
											\end{enumerate}
											\begin{exampleblock}{Answer: B}\end{exampleblock}
											\begin{block}{Explanation}
												Credit risk often monthly.
											\end{enumerate}
										\end{block}
									\end{frame}
									
									\begin{frame}{Q283 --- ALM Backtesting}
										\begin{block}{Question}
											Asset-liability management models may be tested:
											\begin{enumerate}
												\item Daily
												\item Quarterly
												\item Annually
												\item Never
											\end{enumerate}
											\begin{exampleblock}{Answer: B}\end{exampleblock}
											\begin{block}{Explanation}
												ALM models typically quarterly.
											\end{block}
										\end{frame}
										
										\begin{frame}{Q284 --- Validation Independence}
											\begin{block}{Question}
												Validation must be:
												\begin{enumerate}
													\item Done by developers
													\item Independent from development
													\item Optional
													\item Only external
												\end{enumerate}
											\end{block}
											\begin{exampleblock}{Answer: B}\end{exampleblock}
											\begin{block}{Explanation}
												Independence is essential.
											\end{block}
										\end{frame}
										
										\begin{frame}{Q285 --- Validation Report Contents}
											\begin{block}{Question}
												A validation report typically includes:
												\begin{enumerate}
													\item Findings and issues to address
													\item Pass/fail rating
													\item Assumptions reviewed
													\item All of the above
												\end{enumerate}
											\end{block}
											\begin{exampleblock}{Answer: D}\end{exampleblock}
											\begin{block}{Explanation}
												Reports include findings, ratings, assumptions, more.
											\end{block}
										\end{frame}
										
										\begin{frame}{Q286 --- Model Approval}
											\begin{block}{Question}
												Formal approval or rejection of a model is the responsibility of:
												\begin{enumerate}
													\item Developers
													\item Senior management
													\item Only auditors
													\item Only regulators
												\end{enumerate}
											\end{block}
											\begin{exampleblock}{Answer: B}\end{exampleblock}
											\begin{block}{Explanation}
												Senior management approves models.
											\end{block}
										\end{frame}
										
										\begin{frame}{Q287 --- Validation Timing}
											\begin{block}{Question}
												Initial validation occurs:
												\begin{enumerate}
													\item After deployment
													\item During or after design, before implementation
													\item Only at retirement
													\item Never
												\end{enumerate}
											\end{block}
											\begin{exampleblock}{Answer: B}\end{exampleblock}
											\begin{block}{Explanation}
												Initial validation before deployment.
											\end{block}
										\end{frame}
										
										\begin{frame}{Q288 --- Change-Based Validation}
											\begin{block}{Question}
												Change-based validation is triggered by:
												\begin{enumerate}
													\item Calendar schedule
													\item Material change to model
													\item Random selection
													\item Never
												\end{enumerate}
											\end{block}
											\begin{exampleblock}{Answer: B}\end{exampleblock}
											\begin{block}{Explanation}
												Material changes trigger change-based validation.
											\end{block}
										\end{frame}
										
										\begin{frame}{Q289 --- Retirement Documentation}
											\begin{block}{Question}
												Model retirement requires:
												\begin{enumerate}
													\item No documentation
													\item Proper documentation and inventory update
													\item Only verbal confirmation
													\item Only email
												\end{enumerate}
											\end{block}
											\begin{exampleblock}{Answer: B}\end{exampleblock}
											\begin{block}{Explanation}
												Retirement needs documentation and inventory update.
											\end{block}
										\end{frame}
										
										\begin{frame}{Q290 --- GenAI Governance Levers}
											\begin{block}{Question}
												Internal governance levers for GenAI include:
												\begin{enumerate}
													\item Ownership assignment
													\item Risk-based review
													\item Escalation paths
													\item All of the above
												\end{enumerate}
											\end{block}
											\begin{exampleblock}{Answer: D}\end{exampleblock}
											\begin{block}{Explanation}
												All are internal governance levers.
											\end{block}
										\end{frame}
										
										\begin{frame}{Q291 --- GenAI External Governance}
											\begin{block}{Question}
												External governance constraints include:
												\begin{enumerate}
													\item Tracking regulatory changes
													\item Model provider terms
													\item Emerging standards
													\item All of the above
												\end{enumerate}
											\end{block}
											\begin{exampleblock}{Answer: D}\end{exampleblock}
											\begin{block}{Explanation}
												External constraints cover regulation, providers, standards.
											\end{block}
										\end{frame}
										
										\begin{frame}{Q292 --- GenAI Adaptive Practices}
											\begin{block}{Question}
												Adaptive GenAI governance includes:
												\begin{enumerate}
													\item Continuous monitoring
													\item Red-teaming
													\item Scenario planning
													\item All of the above
												\end{enumerate}
											\end{block}
											\begin{exampleblock}{Answer: D}\end{exampleblock}
											\begin{block}{Explanation}
												Adaptive practices: monitoring, red-teaming, scenario planning.
											\end{block}
										\end{frame}
										
										\begin{frame}{Q293 --- GenAI Risk Assessment}
											\begin{block}{Question}
												Strategic GenAI risk assessment asks:
												\begin{enumerate}
													\item Is GenAI the right tool?
													\item What are the risks to safety, legality, reputation?
													\item Both A and B
													\item Neither
												\end{enumerate}
											\end{block}
											\begin{exampleblock}{Answer: C}\end{exampleblock}
											\begin{block}{Explanation}
												Both strategic fit and risk assessment are needed.
											\end{block}
										\end{frame}
										
										\begin{frame}{Q294 --- GenAI Data Governance}
											\begin{block}{Question}
												GenAI data governance should address:
												\begin{enumerate}
													\item Data being used
													\item Data protection policies
													\item Provider input usage
													\item All of the above
												\end{enumerate}
											\end{block}
											\begin{exampleblock}{Answer: D}\end{exampleblock}
											\begin{block}{Explanation}
												All data governance aspects apply.
											\end{block}
										\end{frame}
										
										\begin{frame}{Q295 --- GenAI Technical Controls}
											\begin{block}{Question}
												GenAI technical risk controls include:
												\begin{enumerate}
													\item Guardrails
													\item Human review triggers
													\item Usage limits
													\item All of the above
												\end{enumerate}
											\end{block}
											\begin{exampleblock}{Answer: D}\end{exampleblock}
											\begin{block}{Explanation}
												Guardrails, review triggers, usage limits.
											\end{block}
										\end{frame}
										
										\begin{frame}{Q296 --- GenAI Org Structures}
											\begin{block}{Question}
												GenAI organizational structures should ensure:
												\begin{enumerate}
													\item Clear ownership
													\item Staff training
													\item Aligned incentives
													\item All of the above
												\end{enumerate}
											\end{block}
											\begin{exampleblock}{Answer: D}\end{exampleblock}
											\begin{block}{Explanation}
												Ownership, training, and incentives are key.
											\end{block}
										\end{frame}
										
										\begin{frame}{Q297 --- GenAI External Environment}
											\begin{block}{Question}
												GenAI external environment considerations include:
												\begin{enumerate}
													\item Changing laws
													\item Best practices
													\item Standards alignment
													\item All of the above
												\end{enumerate}
											\end{block}
											\begin{exampleblock}{Answer: D}\end{exampleblock}
											\begin{block}{Explanation}
												External environment = laws, best practices, standards.
											\end{block}
										\end{frame}
										
										\begin{frame}{Q298 --- GenAI Final Thoughts}
											\begin{block}{Question}
												Effective AI/ML validation requires:
												\begin{enumerate}
													\item Qualified validation team
													\item End-to-end review
													\item Balance of performance and other factors
													\item All of the above
												\end{enumerate}
											\end{block}
											\begin{exampleblock}{Answer: D}\end{exampleblock}
											\begin{block}{Explanation}
												All are needed for effective validation.
											\end{block}
										\end{frame}
										
										\begin{frame}{Q299 --- Responsible Governance GenAI}
											\begin{block}{Question}
												Responsible GenAI governance requires:
												\begin{enumerate}
													\item One-time approval only
													\item Adaptive oversight with ongoing monitoring
													\item No oversight
													\item Only IT approval
												\end{enumerate}
											\end{block}
											\begin{exampleblock}{Answer: B}\end{exampleblock}
											\begin{block}{Explanation}
												Ongoing adaptive oversight is essential.
											\end{block}
										\end{frame}
										
										\begin{frame}{Q300 --- Final Question}
											\begin{block}{Question}
												The most important principle across all AI/ML governance is:
												\begin{enumerate}
													\item Maximise accuracy at all costs
													\item Balance performance, fairness, explainability, and governance
													\item Use the most complex models
													\item Ignore regulatory requirements
												\end{enumerate}
											\end{block}
											\begin{exampleblock}{Answer: B}\end{exampleblock}
											\begin{block}{Explanation}
												Effective AI governance balances performance, fairness, explainability, governance.
											\end{block}
										\end{frame}
										
						
										
									\end{document}